\documentclass[11pt]{article}
% Omit random trailer identifiers for reproducible, text-auditable PDF metadata.
\pdftrailerid{}
\usepackage[T1]{fontenc}
\usepackage[utf8]{inputenc}
\usepackage{lmodern}
\usepackage[margin=1in]{geometry}
\usepackage{amsmath,amssymb,amsthm,mathtools}
\usepackage{enumitem,booktabs,array,longtable}
\usepackage{xcolor,microtype}
\usepackage{etoolbox}
\usepackage{xurl}
\usepackage{tikz}
\usetikzlibrary{arrows.meta,positioning}
\usepackage[colorlinks=true,linkcolor=blue!45!black,citecolor=green!35!black,urlcolor=blue!55!black]{hyperref}
\usepackage[capitalize,nameinlink]{cleveref}
\hypersetup{pdftitle={Directed Recovery Structures in Tape Computations},pdfsubject={Local working research manuscript, October 25, 2026},pdfkeywords={matrix transposition, multitape Turing machines, storage codes, directed incidence, index coding}}
\numberwithin{equation}{section}
% ---------------------------------------------------------------------------
% Theorem environments
% ---------------------------------------------------------------------------
\theoremstyle{plain}
\newtheorem{theorem}{Theorem}[section]
\newtheorem{lemma}[theorem]{Lemma}
\newtheorem{corollary}[theorem]{Corollary}
\newtheorem{proposition}[theorem]{Proposition}
\newtheorem{conjecture}[theorem]{Conjecture}
\newtheorem{question}[theorem]{Question}
\newtheorem{problem}[theorem]{Problem}
\theoremstyle{definition}
\newtheorem{definition}[theorem]{Definition}
\newtheorem{remark}[theorem]{Remark}
\newtheorem{example}[theorem]{Example}

% The Lean proof is the proof; the text gives a faithful outline of it.
\newenvironment{outline}[1][Outline of the formal proof]{\begin{proof}[#1]}{\end{proof}}
% Results that are not formalized get complete proofs, labelled as such.
\newenvironment{informalproof}[1][Proof (not formalized)]{\begin{proof}[#1]}{\end{proof}}

% ---------------------------------------------------------------------------
% Lean references
%   \lean{name}            a declaration, relative to the namespace DirectedRecovery
%   \formal{file}{names}   the formal counterpart of a numbered result: a file relative to
%                          lean/DirectedRecovery/ and a comma-separated list of declarations
%   \partformal{file}{names}{gap}
%                          a numbered result whose formal counterpart covers all but the
%                          part described in {gap}; the gap is proved or cited in the text
%   \notformalized         a numbered result that has no formal counterpart
%   \cited{where}          a result of an earlier paper used without proof: where it is proved,
%                          and whether it is Lean-verified there
% ---------------------------------------------------------------------------
% Lean names are typeset like paths, with line breaks allowed after "_", "." and "/".
\DeclareUrlCommand\leanname{\def\UrlFont{\ttfamily\upshape}}
\newcommand{\lean}[1]{\leanname{#1}}
\newcommand{\leanfile}[1]{\leanname{#1}}
\newcommand{\leanlist}[1]{\begingroup\def\leansep{}%
  \renewcommand*{\do}[1]{\leansep\leanname{##1}\def\leansep{, }}\docsvlist{#1}\endgroup}
\newcommand{\formal}[2]{\par\noindent{\small\upshape\raggedright\textsf{Lean:} \leanlist{#2}\unskip
  \nobreak\hfil\penalty50\hskip1em\hbox{}\nobreak\hfill[\leanfile{#1}]\parfillskip=0pt\par}%
  \nobreak\smallskip\nobreak\noindent\ignorespaces}
\newcommand{\partformal}[3]{\par\noindent{\small\upshape\raggedright\textsf{Lean (partial):} \leanlist{#2}\unskip
  \nobreak\hfil\penalty50\hskip1em\hbox{}\nobreak\hfill[\leanfile{#1}]\parfillskip=0pt\par}%
  \par\noindent{\small\upshape\textsf{Not formalized:} #3}\par\nobreak\smallskip\nobreak\noindent\ignorespaces}
\newcommand{\notformalized}{\par\noindent{\small\upshape\textsf{Not formalized;} a complete proof is given
  below.}\par\nobreak\smallskip\nobreak\noindent\ignorespaces}
\newcommand{\cited}[1]{\par\noindent{\small\upshape\textsf{Cited:} #1}\par\nobreak\smallskip\nobreak\noindent\ignorespaces}

% ---------------------------------------------------------------------------
% Notation
% ---------------------------------------------------------------------------
\newcommand{\N}{\mathbb{N}}
\newcommand{\Z}{\mathbb{Z}}
\newcommand{\Q}{\mathbb{Q}}
\newcommand{\R}{\mathbb{R}}
\newcommand{\C}{\mathbb{C}}
\newcommand{\F}{\mathbb{F}}
\newcommand{\T}{\mathsf{T}}
\newcommand{\clog}[1]{\lceil \log_2 #1 \rceil}
\newcommand{\abs}[1]{\lvert #1 \rvert}
\newcommand{\card}[1]{\lvert #1 \rvert}
\newcommand{\eps}{\varepsilon}
\newcommand{\prev}{\mathrm{prev}}
\newcommand{\cost}{\mathrm{cost}}
\newcommand{\gn}{\mathrm{gn}}
\newcommand{\Sh}{\mathrm{Sh}}
\newcommand{\cl}{\mathrm{cl}}
\newcommand{\Aut}{\mathrm{Aut}}
\newcommand{\minrk}{\mathrm{minrk}}
\newcommand{\MAIS}{\mathrm{MAIS}}
\newcommand{\Crux}{\mathrm{Crux}}
\newcommand{\TCP}{\mathrm{TCP}}
\DeclareMathOperator{\rank}{rank}
\DeclareMathOperator{\dist}{dist}
\DeclareMathOperator{\poly}{poly}

\DeclareMathOperator{\tr}{tr}
\DeclareMathOperator{\supp}{supp}
\newcommand{\norm}[1]{\lVert#1\rVert}

\setlist{itemsep=3pt,topsep=5pt}
\setlength{\emergencystretch}{2em}
\title{Directed Recovery Structures in Tape Computations}
\author{Hans Brende \and GPT-6 (OpenAI)}
\date{October 2026 (Draft)}

\begin{document}
\maketitle
\begin{abstract}
Depth reduction turns an oblivious tape computation of matrix transposition into an index code whose side information consists of bounded dependency cones. For fixed-depth cones of linear-time native schedules, we prove that, after paid degree trimming, only $O(L/\log d)$ matched input and demand ports lie within any fixed incidence distance. Here $L=d^2b$ is the bit length of a $d\times d$ matrix. The theorem uses the full syntactic dependency graph and does not assume that the computation is correct.

We obtain half-rate bounds for bounded complex linear recovery, for arbitrary finite-alphabet codes on uniformly hyperfinite incidence graphs, and for directed cyclic components with at most two internal inputs. Cheap native executions show that neither hyperfiniteness nor sublinear reduction to those cyclic components is universal. Other native examples retain many connected demand pairs while admitting an all-code deficit independent of fixed cone depth; a reusable row-word theorem isolates the interface behind that deficit.

An algebraic appendix gives matrix-label certificates uniform in vector width and applies them to an expanding characteristic-two family. A central-layer argument transfers this bound to nonlinear codes with finite-$2$-group endomorphic recovery expressions, including a quantitative payment for exceptional rows.

Finite counting and algebraic interfaces are verified in Lean. The native bridge and its applications are partly formalized; the analytic, cyclic-component, group-transfer, and graph-construction arguments have written proofs with their status stated individually. The general finite-alphabet coding question, unrestricted cheap depth reduction, and a paid treatment of adaptive schedules remain open. No unrestricted transposition time bound follows here.

\end{abstract}

\begin{quote}\small
\textbf{Contributions.} Hans Brende chose the problem and directed the research.
Codex, a GPT-6-based assistant developed by OpenAI, developed the arguments and
prepared this manuscript under his direction. The initial draft contains written
proofs; Lean verification is recorded result by result as it is completed.
\end{quote}

\clearpage
\tableofcontents
\clearpage
\section{Introduction}\label{sec:introduction}

A tape machine can compute a local linear transformation cheaply without making the corresponding recovery graph simple. Conversely, a complicated recovery graph need not support a high-rate code. These two observations create an applicability problem for coding approaches to matrix transposition: one needs a structural condition that follows from the actual execution, together with a coding theorem that uses that condition without silently strengthening it.

This paper studies that interface. Its central object is a directed recovery relation, represented by a bipartite graph with two copies of every coordinate. The left copy is an input port; the right copy is the receiver demanding that coordinate. An edge records available side information. The condition supplied by tape geometry is that almost all matched ports are far apart. Ordinary bipartiteness of the incidence graph is automatic. Symmetry of the recovery relation, local bipartiteness after symmetrization, cheap separators, and disconnected matched ports are additional assertions, and some fail for inexpensive native executions.

\subsection{The transposition problem and the coding interface}
Fix a deterministic multitape Turing machine with a finite state set, a fixed finite work alphabet, a read-only two-way input tape, finitely many two-way work tapes, and an append-only binary output. The input consists of a $d\times d$ matrix written in row-major order, with
\[
 b(d)=\max\{1,\lceil\log_2 d\rceil\},\qquad L=d^2b(d).
\]
The machine must output the transposed matrix in the same format. Statement~(T) asks whether one such machine can perform this task in $O(L)$ time for all~$d$. The earlier papers develop its relation to integer multiplication, causal coding, restricted machine architectures, and bit-versus-word transposition~\cite{BrendeTapeCoding,BrendeLocalCoding,BrendeBitWord,BrendeCurved}. The native incidence-distance theorem allows every entry width $b\ge1$; the explicit machines below use the canonical width.

An oblivious machine has one head schedule and one set of output times at each input length. A \emph{genie} reveals the whole values of selected computation steps. If deleting dependencies into those steps leaves bounded depth, each output is a function of the genie and a bounded set of raw input coordinates. If the machine transposes correctly, fixing the genie leaves a storage code for these side-information sets. A code-size deficit larger than the genie cost gives a contradiction. This is the elementary decomposition of~\cite[Theorem~D]{BrendeCurved}; we recall the required finite counting argument in \cref{sec:preliminaries}.

The unresolved issue is not the counting argument. It is the structure and capacity of the full cone relations. A lower bound for linear codes alone does not control an arbitrary genie, and a graph theorem does not apply until its graph hypothesis has been justified from the execution.

\subsection{Results and their scope}
The main results separate structural information supplied by a trace from the additional hypotheses needed by each coding theorem.

\begin{enumerate}[label=\textup{(\arabic*)}]
\item \textbf{A structural theorem for actual cones.} For every fixed depth and degree bound, a linear-time native schedule has only $O(L/\log d)$ coordinates whose matched ports are within any fixed incidence distance. Conditioning on heavily used inputs makes both incidence degrees bounded at an arbitrarily small fixed coordinate cost. The proof realizes short incidence paths as bounded-congestion geometric transport along alternating forward and reverse copies of the \emph{same} head schedule. Reversing geometric paths does not invert the computation's local functions.

\item \textbf{A directed analytic coding theorem.} Bounded-degree incidence systems satisfying this paired-distance condition have complex linear storage rate at most $1/2+o(1)$ when block size and recovery coefficients are uniformly bounded. Independent phases on incidence components in a local limit replace the symmetric graph hypothesis used in the fourth paper. The argument does not establish a finite-field or nonlinear theorem.

\item \textbf{An all-code theorem, and the failure of its universal structural premise.} Uniform hyperfiniteness together with paired-distance separation gives a half-rate bound for every finite alphabet and every recovery function. A native expander-shear construction shows that actual cone incidence need not be hyperfinite after cheap genie deletion. Its own cone relations nevertheless satisfy an all-code half bound, by a different argument using physical row lanes.

\item \textbf{A directed nonlinear bound with a measured cleanup cost.} If every strongly connected component has internal indegree at most two, all codes obey half rate, apart from an explicit bidirected-triangle correction. Native transport bounds that correction after any eligible deletion. A fixed linear-time three-bank XOR machine nevertheless requires a linear deletion before entering this class. Thus cheap depth reduction alone cannot supply this normalization.

\item \textbf{Connected demands with a positive coding deficit.} A conjugated native construction keeps a positive fraction of matched ports connected after the prescribed deletions. A general theorem turns a bounded cover by monotone row words into an independent rectangle and an all-code deficit. A full-cone analysis verifies this cover for the example. The resulting nonlinear rate deficit is independent of the fixed depth parameter. It is a theorem about these relations and an explicitly stated row interface, not a normal form for arbitrary tape computations.
\end{enumerate}

The algebraic appendix supplies characteristic-two half-rate bounds at arbitrary vector width on a separate expanding family. A central-layer theorem identifies the exact rate benchmark for recovery expressions over finite $p$-groups with the all-width vector-linear benchmark on the same relation. On the expanding family it therefore bounds a specified nonlinear class, uniformly in group size, and pays explicitly for exceptional rows. It does not include arbitrary nonlinear maps or realize this family by a tape schedule.

The principal structural statement is \cref{thm:nb-main}. The following table distinguishes the coding conclusions and their missing premises. Each deficit is relative to the full coordinate count; stated exceptional coordinates must be paid for.

\begin{table}[htbp]
\centering\small
\begin{tabular}{@{}>{\raggedright\arraybackslash}p{3.2cm}>{\raggedright\arraybackslash}p{3.1cm}>{\raggedright\arraybackslash}p{2.5cm}>{\raggedright\arraybackslash}p{5.6cm}@{}}
\toprule
Relation hypothesis & Code class & Conclusion & Additional native obligation \\
\midrule
Paired distance and uniform hyperfiniteness & Every finite alphabet & Half rate asymptotically & Supply small separators; false for some cheap native traces. \\[4pt]
At most two internal inputs per cyclic component & Every finite alphabet & Half rate plus triangle and deletion costs & Supply small cyclic cleanup; the three-bank example requires linear cost. \\[4pt]
Bounded row-word cover, sublinear exceptions & Every finite alphabet and every induced code & Fixed positive deficit & Extract the cover from the full trace; verified here for one example. \\[4pt]
Bounded degree and paired distance & Complex linear, fixed width and bounded coefficients & Half rate asymptotically & Show actual genie fibres belong to this code class. \\[4pt]
The constructed binary matrix-label family & All-width characteristic-two linear; finite-$2$-group expression codes & Half rate plus $\Delta/(2m)$ & Supply the labels and code language on an actual relation. \\
\bottomrule
\end{tabular}
\caption{Theorems control different interfaces. None supplies every premise of an unrestricted native lower bound.}
\label{tab:regimes}
\end{table}

\subsection{Relation to earlier results}
The first paper's causal block-network formulation and common-monotone-pass bounds remain distinct from the step-level incidence approach~\cite{BrendeTapeCoding}. The second paper shows why data-independent geometric conditions and cheap encoders alone do not force the desired coding restrictions~\cite{BrendeLocalCoding}. The bit-to-word reduction of the third paper concerns the computational task rather than the graph hypothesis studied here~\cite{BrendeBitWord}. The fourth paper supplies the depth-to-index-code decomposition, home-based depth reduction, and a symmetric analytic half-rate theorem~\cite{BrendeCurved}. Its accompanying memo extends the weighted all-code circular-clique range from $14/5$ to every ratio below $3$~\cite{BrendeCircularMemo}; no homomorphism from arbitrary native cones to those targets is assumed here.

Storage codes, guessing games, and index coding have an established duality theory~\cite{Mazumdar,ArbabjolfaeiKim}. We use its elementary fibre implication, proved below, and do not need an asymptotic duality theorem to convert a storage bound into a genie lower bound. The analytic proof uses the local-limit and finite-trace framework for unimodular networks~\cite{AldousLyons}. The new issue is the directed incidence geometry and its connection to actual input and output ports.

Two distinctions govern the conclusions. First, the native constructions are inexpensive executors, not transposition algorithms; their role is to test proposed universal graph properties. Second, although their coding bounds exclude transposition on the particular schedules, that algorithmic exclusion is already available from the first paper's pass machinery. The contribution here is the all-code statement for the full cone relation, uniformly over the specified genies and coordinate restrictions.

\subsection{Quantifier conventions}
Unless otherwise stated, machine resources, time constants, cone depth, incidence degree, and vector width are fixed before $d$ or the number of coordinates tends to infinity. An error $o_D(L)$ may have a convergence threshold depending on the fixed depth~$D$. A rate deficit described as independent of~$D$ has one positive leading constant valid for every fixed~$D$; this does not make its error bound uniform when $D$ grows. The all-code results quantify over every finite alphabet and every storage code on the stated relation. The group-expression results instead quantify over their explicitly restricted alphabets and recovery language; group size and expression length may grow. No numerical experiment is a premise of a theorem.

\paragraph{Reading the proofs.} A result marked \textsf{Lean} is proved by the cited Lean declaration; its English outline is a convenience, and the formal statement is authoritative. Finite conditioning, path counting, row-word deficits, and matrix certificates supply reusable verified interfaces. The native structural bridge, its row-word application, and two algebraic applications have explicitly stated formalization gaps. The analytic, cyclic-component, group-expression, and remaining native arguments have complete written proofs. Imported results identify their source and verification scope. \Cref{sec:verification} and the declaration index give the detailed status.

The text develops native separation first, then coding bounds and their limits, followed by native examples. The row-word theorem is stated independently before its machine application. Long rotation and phase-verification proofs are collected in \cref{sec:native-details}; algebraic certificates, group transfer, and the torsion diagnostic are in \cref{sec:mc-section,sec:torsion}.

\section{Recovery relations and the cost of conditioning}\label{sec:preliminaries}

Write $[m]=\{0,\ldots,m-1\}$. For a matrix of entry width~$b$, identify its raw coordinates with
\[
 V=[d]\times[d]\times[b],\qquad
 \operatorname{addr}(i,j,s)=(id+j)b+s,
 \qquad \Pi(i,j,s)=(j,i,s).
\]
Thus $\Pi$ is an involution, and the raw bit at~$p$ is demanded at physical output address $\operatorname{addr}(\Pi(p))$. In particular, for $p=(i,j,s)$ its raw row is~$i$ and its physical output row is~$j$. These two row coordinates must remain distinct in arguments about actual tape geometry. We sometimes write $\Pi(p)$ for the corresponding output address when the address identification is understood.

\begin{definition}[Storage code]\label{def:pre-storage}
Let $V$ be a finite coordinate set and $N(q)\subseteq V$ a recovery set for every $q\in V$. A storage code over a finite alphabet $\mathcal A$ is a set $\mathcal C\subseteq\mathcal A^V$ such that
\[
 x,y\in\mathcal C,\quad x|_{N(q)}=y|_{N(q)}
 \quad\Longrightarrow\quad x_q=y_q.
\]
Equivalently, each coordinate is a function of its recovery set on~$\mathcal C$. The functions may be nonlinear and may depend on the code. For $a=|\mathcal A|\ge2$, the rate of a nonempty code is $|V|^{-1}\log_a|\mathcal C|$.
\end{definition}

The induced relation on $Q\subseteq V$ is $N_Q(q)=N(q)\cap Q$ for $q\in Q$. A code on this induced relation need not extend to one on~$V$. Conversely, projecting an arbitrary code on~$V$ to~$Q$ need not give a code for $N_Q$: the omitted coordinates must first be fixed.

\begin{lemma}[Paid conditioning]\label{lem:pre-conditioning}
\notformalized
Let $H=V\setminus Q$. Every nonempty fibre of a storage code~$\mathcal C$ obtained by fixing the coordinates in~$H$, and then restricting to~$Q$, is a storage code for~$N_Q$. If every code for~$N_Q$ has at most $a^u$ words, then
\[
 |\mathcal C|\le a^{|H|+u}.
\]
\end{lemma}
\begin{informalproof}
Two words in a fixed fibre agree on~$H$. Agreement on $N(q)\cap Q$ therefore gives agreement on all of~$N(q)$, so it determines~$q$. There are at most $a^{|H|}$ fibres, each of size at most~$a^u$.
\end{informalproof}

\begin{definition}[Paired incidence]\label{def:pre-incidence}
The incidence graph $E(N)$ has vertex set $\{L_p:p\in V\}\sqcup\{R_q:q\in V\}$, with an edge $L_pR_q$ exactly when $p\in N(q)$. The pair $L_p,R_p$ is the \emph{demand pair} for~$p$; it is not an extra graph edge. Let
\[
 b_E(r)=\bigl|\{p\in V:\dist_E(L_p,R_p)\le r\}\bigr|,
\]
where disconnected vertices have distance~$\infty$. A sequence on $|V|\to\infty$ satisfies the \emph{paired-distance condition} if $b_E(r)=o(|V|)$ for every fixed integer~$r$.
\end{definition}

The two incidence degrees encode different restrictions: $\deg(R_q)=|N(q)|$ bounds recovery-set size, whereas $\deg(L_p)$ bounds the reuse of a coordinate across recovery sets. Bounded cone depth immediately gives the former, not the latter. If $N$ is symmetric, the incidence graph is the bipartite double cover of its coordinate graph. In that special case a path from $L_p$ to $R_p$ projects to an odd closed walk based at~$p$. For directed recovery, forgetting which endpoint supplies information loses this interpretation.

\begin{definition}[Index code]
For binary inputs on~$V$, an index code of length~$m$ for~$N$ is a map $G:\{0,1\}^V\to\{0,1\}^m$ such that $G(x)$ together with $x|_{N(q)}$ determines~$x_q$ for every~$q$.
\end{definition}

\begin{lemma}[Fibre bound]\label{lem:pre-fibre}
\notformalized
Every fibre of an index code for~$N$ is a binary storage code for~$N$. Consequently, if all such storage codes have size at most $2^{|V|-s}$, every index code has length at least~$s$.
\end{lemma}
\begin{informalproof}
Inside a fibre the broadcast is fixed, so agreement on the recovery set determines the demanded bit. The at most $2^m$ fibres partition $2^{|V|}$ inputs. Hence $2^{|V|}\le2^m2^{|V|-s}$.
\end{informalproof}

For a native computation, let $I_\Gamma(\ell)$ be the raw input coordinates in the dependency cone of physical output~$\ell$ after a genie~$\Gamma$ is fixed. The relation relevant to transposition is
\begin{equation}\label{eq:transpose-relation}
 N_\Gamma(q)=I_\Gamma(\Pi(q)).
\end{equation}
If the machine correctly transposes every input and the schedule is oblivious, its revealed genie word is an index code for this one relation. If the whole-step value width is~$\beta$, that word has at most $\beta|\Gamma|$ bits. Thus a storage bound of $2^{(1-\rho)L+o(L)}$ contradicts a genie of cost at most~$\eps L$ whenever $\eps<\rho$. The schedule may be studied without any correctness assumption; correctness is used only in this last identification of the demanded output with the raw bit~$q$.

All entropies below are Shannon entropies in bits unless a base is displayed. For a random word supported on a storage code, every recovery relation holds almost surely. Deterministic projection arguments therefore bound arbitrary such distributions as well as uniform distributions on the code. We state cardinality bounds explicitly so that no support-size conclusion is inferred merely from a low entropy under a particular input law.

\section{Paired distances in actual tape cones}\label{sec:nb-section}

The recovery relation extracted from a computation is directed: an input
coordinate may occur in the cone of an output without the converse
occurrence.  The appropriate geometric object retains separate copies of
the input and output ports.  We prove that, after a paid bound on input
reuse, almost every matched pair of ports is far apart in this incidence
graph.  No assumption that the computation transposes is needed.

\subsection{The computation DAG and the incidence graph}

Fix integers $d\ge2$ and $b\ge1$, and write
\[
 L=d^2b,\qquad W=db,\qquad [L]=\{0,\ldots,L-1\}.
\]
Identify a coordinate $(i,j,s)$, with $0\le i,j<d$ and $0\le s<b$, with
the raw address $(id+j)b+s$.  The transpose permutation is
\begin{equation}\label{eq:nb-transpose}
 \Pi\bigl((id+j)b+s\bigr)=(jd+i)b+s.
\end{equation}
Throughout this section, logarithms and entropies have base two.

Consider an execution with one read-only input tape, $k\ge1$ work tapes,
and exactly $L$ output appends in $T$ steps.  Each head starts at zero and
moves by at most one cell per step.  At most one bit is appended per step.
Write $x_0(t)$ for the input-head position and $x_j(t)$ for work-head
position $j$ immediately before step $t$.  A work cell is written at every
visit, including a visit which preserves its symbol.  Let
\[
 \operatorname{prev}_j(t)=\max\{u<t:x_j(u)=x_j(t)\}
\]
when this set is nonempty.  We first take the schedule to be oblivious:
these positions, the append times, and $T$ depend on the input length
alone.  The structural argument will also apply to a single frozen
adaptive trace.

\begin{definition}[Whole-step cones]\label{def:nb-cones}
The value $v_t$ of step $t$ consists of the state after that step, all
symbols written on the $k$ work tapes, and the appended output bit or
a symbol indicating no append.  Fix a constant integer $\beta$ large
enough to encode one such value.  A genie $\Gamma\subseteq[0,T)$
reveals $(v_t)_{t\in\Gamma}$ and has cost $\beta|\Gamma|$ bits.

The computation DAG $H$ has vertices $0,\ldots,T-1$, state edges
$t\to t-1$ for $1\le t<T$, and cell edges
$t\to\operatorname{prev}_j(t)$ whenever the
displayed predecessor exists.  Parallel edges may be identified.  All
edges point backward in time.

A path $t_0\to\cdots\to t_m$ is $\Gamma$-free if
$t_1,\ldots,t_m\notin\Gamma$; its first vertex may belong to $\Gamma$.
The genie has \emph{full depth at most $D$} if every such path has at most
$D$ edges.  Let $C_\Gamma(t)$ be the set of endpoints of $\Gamma$-free
paths from $t$.  If output address $\ell$ is appended at step $t_\ell$,
put
\[
 I_\Gamma(\ell)=\{x_0(u):u\in C_\Gamma(t_\ell)\}\cap[L],
 \qquad N(q)=I_\Gamma(\Pi(q)).
\]
\end{definition}

These are the whole-step conventions of
\cite[the dependency DAG and cone lemma]{BrendeCurved}.  In particular,
a length-zero path is permitted, even when its vertex is revealed.
Here $N(q)$ is a syntactic input set; its definition does not require
that output $\Pi(q)$ equal input $q$.

Each vertex of $H$ has both in-degree and out-degree at most $k+1$.
For the in-degree assertion, a given visit has at most one next visit to
the same cell on each work tape, and at most one next state edge.
Consequently, if the genie has full depth at most $D$, then
\begin{equation}\label{eq:nb-cone-bound}
 |N(q)|\le B,\qquad B=\sum_{h=0}^{D}(k+1)^h.
\end{equation}

\begin{definition}[Paired incidence graph]\label{def:nb-incidence}
The graph $E$ has distinct vertices $L_p,R_p$ for every $p\in[L]$, with
an undirected edge $L_pR_q$ precisely when $p\in N(q)$.  The physical
port represented by $L_p$ is input address $p$; that represented by $R_q$
is output address $\Pi(q)$.  For $Q\subseteq[L]$, let $E_Q$ be the
induced graph on $\{L_p,R_p:p\in Q\}$.  Distances between vertices in
different components are infinite.
\end{definition}

The graph $E$ is an auxiliary bipartite graph on physical ports.  It is
neither the computation DAG $H$ nor the undirected graph obtained by
forgetting the directions of the recovery relation on $[L]$.  Its
bipartiteness is automatic; the assertion below concerns the distances
between the two copies of the \emph{same} coordinate.

\begin{theorem}[Paired-distance bound]\label{thm:nb-main}
\notformalized
Fix $k\ge1$, $D\ge0$, $c>0$, a degree bound $\Delta$, and an integer
$r\ge1$.  Suppose $T\le c(L+1)$ and $\Gamma$ has full depth at most
$D$.  For every $Q\subseteq[L]$ such that $E_Q$ has maximum degree at
most $\Delta$,
\begin{equation}\label{eq:nb-main-bound}
 \bigl|\{p\in Q:\operatorname{dist}_{E_Q}(L_p,R_p)\le r\}\bigr|
   =O_{k,D,c,\Delta,r}\!\left(\frac{L}{\log_2d}\right).
\end{equation}
The constant is uniform over $b\ge1$, the schedule, the genie, and $Q$.
Neither correctness as a transposer nor a bound on $|\Gamma|$ is assumed.
\end{theorem}

The proof uses a geometric relaxation in which selected paths may share
a location and subsequently separate.  These paths are auxiliary objects;
they do not assert that a coded tape symbol carries independent source
bits.  A row-entropy potential controls partial transpose transport in
this relaxation.  Alternating forward and reverse realizations of cone
matchings then converts short paired paths in $E_Q$ into that transport.

\subsection{A partial-transport inequality}

We specify the geometric relaxation so that its use does not import a
routing assumption about the original machine.  There are $K$ tapes,
one unit-speed head on each, and $\rho$ registers without spatial
addresses.  A geometric tick first permits arbitrary transfers among
the current head cells and registers, leaving other locations unchanged,
and then moves each head by at most one cell.  Each selected source
label follows a single trajectory.  Different trajectories may occupy
the same slot and later separate.  The \emph{occupancy} is the number
of selected trajectories at a slot, without regard to encoded values.
We bound occupancy at each placement before and after a local transfer.
Moving heads does not change the locations of the selected labels.

\begin{lemma}[Partial transpose transport]\label{lem:nb-transport}
\notformalized
Let $S\subseteq[L]$ have $m>0$ elements.  Suppose geometric trajectories
start at the distinct input cells $p\in S$, with no selected labels
elsewhere, and end at the distinct cells $\Pi(p)$ of one output tape.
Assume all heads start at zero, every slot has occupancy at most $C$,
and the trajectories use $U$ geometric ticks.  Then
\begin{equation}\label{eq:nb-transport-sharp}
 m\log_2d-m\log_2\!\left(\frac{L+2W}{m}\right)-m
       \le CK(K+\rho+3)U.
\end{equation}
In particular, including $m=0$,
\begin{equation}\label{eq:nb-transport-simple}
 m\log_2d\le CK(K+\rho+3)U+3L.
\end{equation}
\end{lemma}

\begin{informalproof}
Fix all trajectories before choosing any partition.  For each
$\theta\in\{0,\ldots,W-1\}$, partition every tape into the bins
$\lfloor(z-\theta)/W\rfloor$, retaining the tape name as part of the
bin name.  Merge the $K$ bins currently containing heads, together
with all registers, into one active class.  The remaining bins are
separate classes.  The active class contains at most $KW+\rho$ slots.

Choose a source label $V$ uniformly from $S$, and let
$A=\lfloor V/W\rfloor$ be its original input row.  If
$Z_{t,\theta}$ is the class containing its trajectory at boundary $t$,
define
\[
 \Phi_t=m\,\mathbb E_\theta\mathrm H(A\mid Z_{t,\theta}).
\]
This is entropy of a random \emph{source label}, not entropy of the
machine input or the contents of a cell.  For any two class variables
$Z,Z'$, the chain rule gives
\begin{equation}\label{eq:nb-chain-rule}
 \mathrm H(A\mid Z')-\mathrm H(A\mid Z)
       \le \mathrm H(Z\mid Z').
\end{equation}
Thus a new class containing $n$ labels contributes at most $n$ times
the entropy of their old class names to the increase of the unnormalized
potential.

A local transfer changes no label's class: both of its endpoints lie
in the active class.  Analyze a simultaneous movement of the heads as
$K$ successive movements, with all labels held at their physical slots.
A unit movement of a particular head crosses a bin boundary for exactly
one offset $\theta$.  At that offset, the new active class has at most
two old class names, namely the old active class and the newly joined
bin.  A released bin has only the old active name.  All other classes
are unchanged.  By occupancy, the new active class contains at most
$C(KW+\rho)\le C(K+\rho+3)W$ selected labels.  Equation
\eqref{eq:nb-chain-rule} therefore bounds the increase, after averaging
over $\theta$, by $C(K+\rho+3)$.  A pause has zero cost.  Summing gives
\begin{equation}\label{eq:nb-potential-motion}
 \Phi_U-\Phi_0\le CK(K+\rho+3)U.
\end{equation}

Initially, a shifted input bin meets at most two aligned input rows.
Only the input tape contains selected labels, and merging headed bins
merges no additional occupied input bin.  Hence $\Phi_0\le m$.

For the endpoint estimate, first keep all output bins separate.  At
output position $v=\Pi(p)$ the original source row is
$\lfloor(v\bmod W)/b\rfloor$.  This label is $W$-periodic, with exactly
$b$ occurrences of each row in a period.  Thus a shifted $W$-cell bin
containing $n$ selected labels has row multiplicities $n_a\le b$, and
\[
 n\,\mathrm H(A\mid\text{this bin})
     =n\log_2 n-\sum_a n_a\log_2n_a
     \ge n\log_2 n-n\log_2b.
\]
Use $0\log_20=0$.  At most $L/W+2$ shifted bins meet $[L]$.
Convexity of $x\log_2x$ consequently gives
\[
 m\,\mathrm H(A\mid\text{separate output bins})
   \ge m\log_2\!\left(\frac{m}{b(L/W+2)}\right)
   =m\log_2d-m\log_2\!\left(\frac{L+2W}{m}\right).
\]
Merging headed bins is a coarsening of the partition, so it cannot
decrease this conditional entropy.  This is a lower bound for $\Phi_U$.
Together with \eqref{eq:nb-potential-motion} and $\Phi_0\le m$, it proves
\eqref{eq:nb-transport-sharp}.

For $a>0$ and $0<x\le a$,
$x\log_2(a/x)\le a/(e\ln2)$.  Since $L+2W\le2L$ and $m\le L$,
\[
 m+m\log_2\!\left(\frac{L+2W}{m}\right)
 \le \left(1+\frac{2}{e\ln2}\right)L<3L.
\]
Rearranging proves \eqref{eq:nb-transport-simple}.
\end{informalproof}

The factor $b$ in the endpoint estimate is essential.  The statement
holds for arbitrary selected source subsets and every entry width; no
conversion of a word instance into a bit instance is used.

\subsection{Realizing and reversing incidence matchings}

\begin{lemma}[A matching of cone incidences]\label{lem:nb-matching}
\notformalized
Let $\mathcal M$ be a matching in $E$.  There is a geometric realization
of its edges, carrying a label from input cell $p$ to output cell
$\Pi(q)$ for each $L_pR_q\in\mathcal M$, with $K=k+2$ tapes, one
register, occupancy at most $B+1$, and $T$ ticks.  All such realizations
use the same head configurations $h_0$ at entry and $h_T$ at exit.
The geometric paths can be reversed in $2T$ ticks with the same
occupancy, from configuration $h_T$ to $h_0$.
\end{lemma}

\begin{informalproof}
For every edge $L_pR_q$ choose one $\Gamma$-free backward path from
the output step $t_{\Pi(q)}$ to a step reading raw input $p$.  It has
at most $D$ edges.  Any vertex of the computation DAG belongs to at
most $B$ chosen paths: the distinct output steps of paths using that
vertex lie within its forward radius $D$, which contains at most $B$
vertices.  This counts output endpoints, not the number of possible
combinatorial paths.

For the geometric realization, regard the input tape and the
append-output tape as two additional work-capable tapes.  The input
head follows its original trajectory.  The output head occupies the
next append address and advances exactly when an append occurs.
The additional register represents a whole step's record for purposes
of the geometric interface; it is not restricted to the original
control-state component.

Read each chosen DAG path in forward chronological order.  Its label
waits at its raw input cell until the chosen input read.  A state edge
is realized by waiting in the record register from one step to the
next.  A cell edge is realized by waiting in that work cell between
two consecutive visits.  At a path vertex, the incoming head cell or
register and the required outgoing head cell or register are present
at the local interface, so the geometric model permits the transfer.
At the chosen output step the label moves to the current output cell
and remains there.  These operations occur before the original head
movements, and therefore use the original $T$ ticks.

During a work-cell waiting interval all labels there use its same
previous-visit vertex, so there are at most $B$.  Labels in the record
register use one step vertex and satisfy the same bound.  A raw input
cell contains only its own not-yet-departed source label; a completed
output cell contains only its own output label.  The matching supplies
distinct sources and outputs, and the input and output tapes are
distinct.  The bound $B+1$ is thus conservative at every placement.
At the end of the block, all selected labels occupy output ports and
none remains in a work cell or register.

The record register is an augmentation of the geometric model.  A
syntactic state edge of the whole-step DAG may refer to components that
the real control state does not retain.  The argument does not identify
that edge with an independent information channel in the original
machine.  Instead it uses an additional register in precisely the
geometric model covered by \cref{lem:nb-transport}.

To reverse a tick, first preserve all labels and undo the original unit
head movements in one tick.  In a second tick, keep the heads fixed and
reverse the selected local transfers.  Both endpoints of each transfer
are again current head cells or registers.  This reverses the geometric
trajectories, including waiting intervals, with the same occupancy.
It does not invert a transition function or reconstruct a past machine
state.  All matchings use the same original head schedule, including
the two port tapes.  Their forward and reverse configurations therefore
agree exactly at the stated block boundaries.
\end{informalproof}

\begin{lemma}[Disjoint short paired paths]\label{lem:nb-disjoint}
\notformalized
For any fixed odd $s\ge1$, the number $m$ of vertex-disjoint $s$-edge
paths in $E$ whose endpoints are $L_p,R_p$ satisfies
\begin{equation}\label{eq:nb-disjoint-bound}
 m\le A_s\frac{L}{\log_2d},\qquad
 A_s=4sc(B+1)(k+2)(k+6)+3.
\end{equation}
\end{lemma}

\begin{informalproof}
Write $s=2a+1$ and orient each incidence path from its left endpoint
to its right endpoint.  The edges at a fixed position along these
paths form a matching, because the paths are vertex-disjoint.  Realize
the first matching forward, the second backward, and continue
alternating, ending with a forward block.

For example, an incidence path beginning
$L_p,R_v,L_u,R_w$ first carries its label from input $p$ to output
$\Pi(v)$, then from that same output cell back to input $u$, and
then to output $\Pi(w)$.  Every splice identifies the same physical
tape and address.  In particular, $\Pi$ is never performed as a free
operation; it only names which output cell a right port represents.
Forward blocks end at $h_T$, where reverse blocks begin, and reverse
blocks end at $h_0$, where the next forward blocks begin.  No head
repositioning is omitted.  At each junction the selected labels occupy
distinct ports, so occupancy does not accumulate across stages.

The resulting trajectories carry each original input $p$ to output
$\Pi(p)$, starting with every head at zero.  They use $K=k+2$ tapes,
$\rho=1$, occupancy $B+1$, and
\[
 U=(3a+1)T\le2sT\le4scL
\]
ticks.  Apply \eqref{eq:nb-transport-simple} to obtain
\eqref{eq:nb-disjoint-bound}.
\end{informalproof}

\begin{informalproof}[Proof of \cref{thm:nb-main}]
For each $p\in Q$ at paired distance at most $r$, choose one shortest
path in $E_Q$.  Its length is odd.  Put
\[
 V_r=\sum_{j=0}^{r}\max(1,\Delta)^j.
\]
A fixed vertex $z$ belongs to at most $V_r$ chosen paths: the distinct
left endpoints of paths through $z$ are within distance $r$ of $z$,
and there is one chosen path per left endpoint.

For each odd $s\le r$, choose a maximal vertex-disjoint subfamily among
the chosen paths of length $s$.  Every candidate meets a selected path.
A selected path has $s+1$ vertices and therefore meets at most
$(s+1)V_r$ candidates.  \cref{lem:nb-disjoint} gives the explicit bound
\begin{equation}\label{eq:nb-packing-bound}
 \bigl|\{p\in Q:\operatorname{dist}_{E_Q}(L_p,R_p)\le r\}\bigr|
 \le \sum_{\substack{1\le s\le r\\
s\text{ odd}}}
       (s+1)V_r A_s\frac{L}{\log_2d}.
\end{equation}
All parameters in the sum are fixed independently of $b$, the schedule,
and the genie.  This proves the theorem.
\end{informalproof}

\subsection{Input reuse, conditioning, and adaptive traces}

\begin{corollary}[Self-demands]\label{cor:nb-self-demands}
\notformalized
Without any bound on left incidence degrees,
\[
 |\{p\in[L]:p\in N(p)\}|\le A_1\frac{L}{\log_2d}.
\]
\end{corollary}

\begin{informalproof}
The edges $L_pR_p$ for self-demanding coordinates form a matching.
Apply \cref{lem:nb-disjoint} with $s=1$.
\end{informalproof}

\begin{corollary}[Paid degree reduction]\label{cor:nb-trim}
\partformal{IncidenceTrimming.lean}{paid_degree_trimming}{The native cone-size premise and the subsequent paired-distance asymptotic application.}
For an integer $M\ge1$, let
\[
 H_M=\{p\in[L]:\deg_E(L_p)>M\},\qquad Q_M=[L]\setminus H_M.
\]
Then $|H_M|\le BL/M$, and $E_{Q_M}$ has maximum degree at most
$\max(B,M)$.  For each fixed $\delta\in(0,1)$, choosing the fixed
threshold $M\ge B/\delta$ retains at least $(1-\delta)L$ coordinates
and gives \eqref{eq:nb-main-bound} at every fixed radius.
\end{corollary}

\begin{informalproof}
By \eqref{eq:nb-cone-bound},
$\sum_p\deg_E(L_p)=\sum_q|N(q)|\le BL$.
This bounds the number of left degrees exceeding $M$.  Deleting both
ports of each such coordinate cannot increase a degree.  Surviving
left degrees are at most $M$, and all right degrees were at most $B$.
\end{informalproof}

The deletion in \cref{cor:nb-trim} costs a small \emph{fixed}
fraction for a fixed threshold; it is not an uncharged $o(L)$ deletion.
Its coding interpretation is paid conditioning as in
\cref{lem:pre-conditioning}.  \cref{prop:open-transfer} below
spells out the resulting conditional transfer of an all-code half bound.

\begin{remark}[Exactly what is structural]\label{rem:nb-scope}
The premise of \cref{prop:open-transfer} is an additional
coding statement.  \cref{thm:nb-main} alone supplies no code-rate
bound.  It also supplies no genie: the depth $D$ is assumed fixed before
$d$ tends to infinity.  Its constants may grow rapidly with $D$.

Every argument proving \cref{thm:nb-main} applies to one fixed legal
adaptive execution trace with the stated time, output, and depth
properties.  Freeze that trace's head walks and append times and use
them as the geometric schedule.  The only random object in
\cref{lem:nb-transport} is an auxiliary source label, so no distribution
of inputs realizing the trace is required.  This observation does not
produce a common cone relation across inputs, bound the entropy of a
trace, or make the input coordinates independent on a trace fibre.
Those are separate obligations for an adaptive all-input coding
reduction.
\end{remark}

\subsection{Why symmetrizing the recovery relation loses the conclusion}

The distinction between incidence distance and ordinary odd cycles is
necessary even for simple native operations.  The example below forces
directed triangles in every cone relation while respecting the geometric
conditions already proved.

\begin{proposition}[Persistent directed triangles]\label{prop:nb-triangles}
\notformalized
For canonical widths $b=\max(1,\lceil\log_2d\rceil)$, there is a fixed
uniform deterministic oblivious native machine with $O(L)$ runtime
which appends exactly $L$ bits and computes
\[
 Y[j,i,s]=X[j,-j-i\bmod d,s].
\]
For every genie $\Gamma$, its demand-labelled cone relation contains
\begin{equation}\label{eq:nb-triangle-map}
 F(i,j,s)=(j,-i-j\bmod d,s)\in N(i,j,s).
\end{equation}
The map $F$ has $\gcd(3,d)b$ fixed points and otherwise consists of
disjoint directed three-cycles.  Thus the undirected symmetrization of
the recovery relation contains
$(L-\gcd(3,d)b)/3$ vertex-disjoint triangles.  Making it triangle-free
by coordinate deletion requires at least this many deletions.

Moreover, for each fixed $\varepsilon>0$, the machine admits genies of
cost at most $\varepsilon L$ and full depth bounded independently of
$d$.  The same triangle conclusion holds for those genies.
\end{proposition}

\begin{informalproof}
Output rows $j$ in order and, within each row, output entries $i$ in
order.  Obtain the entry directly from input row $j$, column
$c=-j-i\bmod d$.  Read its $b$ bits forward, appending each bit in the
same transition that reads it.  Between successive entries the target
column decreases by one.  An ordinary repositioning costs $O(b)$
input-head steps, and the single wrap within a row costs $O(db)$.
Moving from the final target of one row to the initial target of the
next costs $O(db)$.  The total input motion and emission cost is
therefore $O(d^2b)$.

Here is a fixed-resource implementation of the indexing and motion
control.  First count the raw length $L$ with a little-endian binary
counter, returning its head to zero after each increment.  The total
carry distance is $O(L)$.  Rewind the input head, and recover $d,b$
from the strictly increasing function
$d\mapsto d^2\max(1,\lceil\log_2d\rceil)$ by binary search and
ordinary arithmetic on $O(\log L)$ bits.  This setup takes
$\operatorname{poly}(\log L)=O(L+1)$ time on a fixed number of tapes.
Noncanonical lengths may use an immediate fallback after this test.

At canonical lengths $d\ge2$, every subsequent index and displacement
has $O(b)$ bits.  Per-entry operations are addition, subtraction,
comparison, and modular decrement, implementable in $O(b)$ time by
full scans of a fixed number of binary registers.  A reusable unary
guide of length $b$ controls blocks of $b$ input moves or output
appends.  Each use scans outward and returns, with the input head
paused during the return.  Longer displacements use $O(d)$ such
blocks with $O(b)$ counter work per block; they cost $O(db)$ each.
Hence all entry operations, row wraps, and row transitions still cost
$O(L)$.  The work tapes are registers or guides accessed by monotone
outward-and-return scans from zero, so their walks are home-based.
All motion depends on the public length and counters; payload bits
affect only the emitted values.  This proves uniformity, obliviousness,
and the claimed native cost with a fixed tape alphabet and control.

The step emitting physical output $(j,i,s)$ reads raw input
$(j,-j-i\bmod d,s)$.  The length-zero path at this step is
$\Gamma$-free for every $\Gamma$, including one revealing that step.
Demand relabelling by \eqref{eq:nb-transpose} proves
\eqref{eq:nb-triangle-map}.  Direct calculation gives
\[
 F^2(i,j,s)=(-i-j,i,s),\qquad F^3(i,j,s)=(i,j,s).
\]
Its fixed points satisfy $i=j$ and $3i=0$ in $\mathbb Z/d\mathbb Z$,
so there are exactly $\gcd(3,d)b$ of them.  All other orbits have size
three and give vertex-disjoint triangles after symmetrization.

Finally use the home-based depth-reduction theorem and long-edge
conversion of \cite[Theorem~B and the proof of Corollary~B]{BrendeCurved}.
For fixed machine resources and each $\varepsilon>0$, these give a
genie with cost at most $\varepsilon L$ and constant full depth.
That construction uses the schedule and its runtime, not transpose
correctness.  It therefore applies to this machine.  The triangles
persist because \eqref{eq:nb-triangle-map} held for every genie.
\end{informalproof}

The guaranteed arcs in this example contribute only the incidence
matching $L_{F(q)}R_q$; they do not by themselves connect $L_q$ to
$R_q$ around the directed three-cycle.  There is therefore no conflict
with \cref{thm:nb-main}.  Additional actual cone edges are permitted
throughout the argument.  The example rules out a universal local
bipartiteness conclusion for the symmetrized recovery graph, while
leaving the directed paired-distance condition intact.

\section{Two half-rate bounds for directed recovery}
\label{sec:dc-section}

The distinction between a coordinate and its two incidence ports is essential
here.  A directed recovery relation on a finite set $V$, with $N=|V|$, assigns
a set $\mathcal N(q)\subseteq V\setminus\{q\}$ to each $q\in V$.
Its \emph{paired incidence graph} $E$ has ports $L_p,R_p$, $p\in V$, and
an undirected edge $L_pR_q$ precisely when $p\in\mathcal N(q)$.
Distances below are in this graph; disconnected ports have distance infinity.
Write
\begin{equation}
  b_E(r)=\bigl|\{p\in V:\operatorname{dist}_E(L_p,R_p)\le r\}\bigr|.
  \label{eq:dc-bad-pairs}
\end{equation}
A sequence $(E_n)$ on $2N_n$ ports, with $N_n\to\infty$, has
\emph{paired-distance separation} if $b_{E_n}(r)=o(N_n)$ for every fixed
integer $r\ge1$.  This concerns prescribed pairs of ports, rather than
the girth or connected-component sizes of $E_n$.

We prove a half-rate bound under either of two additional hypotheses.
Uniform hyperfiniteness suffices for every finite-alphabet code.
Without hyperfiniteness, bounded complex linear recovery coefficients
and a fixed block width suffice.  The latter extends the graph hypothesis
of Theorem~A of~\cite{BrendeCurved} from a locally bipartite underlying
graph to paired-distance separation.

\subsection{All codes on hyperfinite incidence graphs}

A storage code over a finite alphabet $\mathcal A$ of size at least two is a set
$\mathcal C\subseteq\mathcal A^V$ such that two codewords agreeing on
$\mathcal N(q)$ agree at $q$, for every $q\in V$.
Equivalently, each coordinate has a recovery function on the neighbour
assignments that occur in the code.  The code need not be linear, and
no distribution or extension property is assumed.

\begin{lemma}[Separated components]
\formal{IncidenceComponents.lean}{storageCode_card_le_of_separated_components}
\label{lem:dc-components}
If $L_p$ and $R_p$ lie in different components of $E$ for every $p\in V$,
then every storage code satisfies
\[
  |\mathcal C|\le |\mathcal A|^{\lfloor N/2\rfloor}.
\]
\end{lemma}
\begin{outline}
Let $\alpha(p)$ and $\beta(p)$ be the components containing $L_p$ and $R_p$.
Choose an order of the components, and record the coordinates $p$ for which
$\alpha(p)$ precedes $\beta(p)$.  This record determines a codeword.
Indeed, process the $\alpha$-components in increasing order.  A coordinate
in the current component was either recorded, or its $\beta$-component is
earlier.  In the latter case every recovery neighbour $u\in\mathcal N(p)$
satisfies $\alpha(u)=\beta(p)$, so its value is already determined.
The equality case $\alpha(p)=\beta(p)$ is excluded by hypothesis.
Compare this order with its reverse. Every coordinate is recorded in exactly
one of the two orders, so one record has at most $\lfloor N/2\rfloor$
coordinates. Its projection is injective on $\mathcal C$.
\end{outline}

\begin{theorem}[Finite separator bound]
\formal{IncidenceSeparator.lean}{storageCode_separator_log_bound}
\label{thm:dc-separator}
Let $Z$ be a set of incidence ports such that each component of $E-Z$
has at most $K$ vertices, where $K\ge1$.  Every nonempty storage code satisfies
\begin{equation}
 \log_{|\mathcal A|}|\mathcal C|
 \le \frac{N+|Z|+b_E(K)}2.
 \label{eq:dc-separator-bound}
\end{equation}
No degree bound is required.
\end{theorem}
\begin{outline}
Let $U$ contain the coordinates with at least one port in $Z$, and also
the coordinates counted by $b_E(K)$.  Thus $|U|\le |Z|+b_E(K)$.
Delete both ports of every coordinate in $U$, obtaining the induced
relation on $Q=V\setminus U$.  Its components still have at most $K$
ports.  If a surviving pair were connected, a shortest path in its
component would have length at most $K-1$, contradicting the choice of $U$.

Fix an assignment to the coordinates in $U$.  The restrictions to $Q$ of
codewords with this assignment form a storage code for the induced relation:
deleted recovery arguments are fixed constants.  \cref{lem:dc-components}
bounds each such fibre by $|\mathcal A|^{|Q|/2}$.  There are at most
$|\mathcal A|^{|U|}$ fibres, giving
$|\mathcal C|\le |\mathcal A|^{|U|+|Q|/2}$ and the claimed inequality.
This conditions on coordinate values; it does not discard side-information
edges without paying for their coordinates.
\end{outline}

\begin{definition}
\label{def:dc-hyperfinite-definition}
\formal{Hyperfinite.lean}{UniformlyHyperfinite}
A sequence $(E_n)$ is \emph{uniformly hyperfinite} if, for each $\eta>0$,
there is an integer $K_\eta$ such that, for every sufficiently large $n$,
one can delete at most $2\eta N_n$ incidence ports and leave components
of size at most $K_\eta$.
\end{definition}

\begin{corollary}[Hyperfinite half-rate bound]
\formal{Hyperfinite.lean}{hyperfinite_rate_limsup_le_half}
\label{cor:dc-hyperfinite}
Suppose $(E_n)$ is uniformly hyperfinite and has paired-distance separation.
For any finite alphabets $\mathcal A_n$ of size at least two, and any
nonempty storage codes $\mathcal C_n$ on these relations,
\[
 \limsup_{n\to\infty}
 \frac{\log_{|\mathcal A_n|}|\mathcal C_n|}{N_n}\le\frac12.
\]
\end{corollary}
\begin{outline}
For fixed $\eta$, \cref{thm:dc-separator} gives a limiting upper bound
$1/2+\eta$, because $K_\eta$ is fixed before $n$ tends to infinity.
Then let $\eta$ decrease to zero.  The finite estimate is uniform in
the alphabet, so it also permits the displayed varying alphabets.
\end{outline}

\subsection{Bounded complex linear recovery}

Let $t\ge1$.  A recovery matrix $A$ has $t\times t$ blocks $A_{qp}$,
with $A_{qq}=0$ and $A_{qp}=0$ unless $L_pR_q\in E$.
Its actual support graph has $L_pR_q$ precisely when $A_{qp}\ne0$;
we may use a larger declared relation $E$.
A linear storage space with these recovery rules is contained in
$\ker(I-A)$.  No symmetry or adjoint relation between $A_{qp}$ and
$A_{pq}$ is imposed.

\begin{theorem}[Directed complex half-rank bound]
\notformalized
\label{thm:dc-complex}
Fix integers $\Delta,t\ge1$ and a real number $B\ge1$.
Let $(E_n)$ have maximum incidence degree at most $\Delta$ and paired-port
separation.  Let $A_n$ be complex recovery matrices of block width $t$
on $E_n$, with every scalar entry of modulus at most $B$.
Then
\begin{equation}
 \liminf_{n\to\infty}
 \frac{\operatorname{rank}_{\mathbb C}(I-A_n)}{tN_n}\ge\frac12.
 \label{eq:dc-complex-rank}
\end{equation}
Consequently every associated complex or real linear storage space has
dimension at most $(1/2+o(1))tN_n$.
\end{theorem}

We give the analytic details because two points matter: a component
marking must preserve unimodularity, and finite nullity passes to the
limit in only one direction.
We use the marked-network framework of Aldous and Lyons~\cite[Sections~2
and~5]{AldousLyons}.  Uniformly rooted finite graphs of bounded degree,
with marks in fixed compact spaces, have locally convergent subsequences;
their limits satisfy the mass-transport identity.  For a unimodular law
$\mu$, bounded measurable operator families equivariant under network
isomorphisms, and independent of the distinguished root, form a von Neumann
algebra.  In its $t\times t$ matrix amplification the normalized trace is
\begin{equation}
 \tau(T)=\mathbb E_\mu\frac1t\operatorname{tr}(T_{oo}),
 \qquad \tau(I)=1.
 \label{eq:dc-trace}
\end{equation}
It is tracial and positive.  These are the same analytic inputs used in
the proof of Theorem~A in~\cite{BrendeCurved}.

\begin{informalproof}[Proof of \cref{thm:dc-complex}]
Form the coordinate graph $G_n$ by joining $p$ and $q$ if either directed
incidence is present.  Its degree is at most $2\Delta$.
Record the two directed incidence flags as discrete marks, and the
corresponding bounded coefficient blocks as compact marks.  In particular,
retain an incidence flag when its coefficient converges to zero.
Choose a subsequence realizing the limsup of the normalized nullities,
then a locally convergent marked subsequence with unimodular limit $\mu$.
The root is uniform among coordinates, and the rooted limit retains the
root's connected coordinate component.  Write $(G,A,o)$ for that limit.
The row and column absolute sums of $A$ are at most $\Delta tB$, so
\[
 \|A\|\le\Delta tB,\qquad \|I-A\|\le c:=1+\Delta tB.
\]

\emph{Separation in the limit.}
For fixed $r,s$, the event that a coordinate within $G_n$-distance $s$
of the root has paired-port distance at most $r$ has probability at most
\[
 \left(\sum_{j=0}^s(2\Delta)^j\right)\frac{b_{E_n}(r)}{N_n}.
\]
It is determined by finitely many discrete incidence flags and therefore
is a clopen local event.  Its limiting probability is zero.  Taking the
countable union over $r,s$ shows that, almost surely, \emph{every}
coordinate $p$ of $G$ has $L_p,R_p$ in different components of the limit
port graph $E$.

\emph{An invariant phase extension.}
Conditionally on this countable graph, give each component of $E$ an
independent uniform phase on the complex unit circle.  This law can be
specified without choosing an invariant enumeration: for a finite list
of ports, give independent phases to its distinct component classes,
and equal phases to ports in the same class.  The cylinder laws are
consistent.  Connectivity of two ports is a countable union of finite-path
events, so their component partition is Borel.  The cylinder integrals,
and then all bounded Borel integrals by a monotone-class argument, define
a Borel probability kernel $\kappa_G$.  Its law is unchanged by any
isomorphism or change of root.

The resulting phase-marked extension $\widehat\mu$ is unimodular.
Indeed, for any nonnegative doubly rooted test function $f$, average
$f$ over $\kappa_G$ to obtain $\bar f$.  The same kernel is used with
either root distinguished.  Applying mass transport to $\bar f$ and
Tonelli's theorem gives
\[
 \mathbb E_{\widehat\mu}\sum_x f(G,o,x)
 =\mathbb E_\mu\sum_x\bar f(G,o,x)
 =\mathbb E_\mu\sum_x\bar f(G,x,o)
 =\mathbb E_{\widehat\mu}\sum_x f(G,x,o).
\]
This construction is a probability-space extension, not an assertion
that component phases are a factor of independent vertex labels.

Let $D_L,D_R$ multiply the $t$ entries at coordinate $p$ by the phase
of the component of $L_p,R_p$, respectively.  They are measurable,
equivariant diagonal unitaries.  An edge $L_pR_q$ equates the phase
at column $p$ with that at row $q$, and hence
\begin{equation}
 AD_L=D_RA.
 \label{eq:dc-intertwining}
\end{equation}
The phases of each matched pair are independent and nonatomic.
By countability they differ at every coordinate simultaneously, almost
surely.  Thus $D_L-D_R$ is injective; no bounded inverse is asserted.

\emph{The tracial half bound.}
Let $K=\ker(I-A)$ and let $P_K$ be its orthogonal projection.
If $y\in K$ and $D_Ly\in K$, then
$D_Ly=AD_Ly=D_RAy=D_Ry$, so $y=0$.
Consequently $T=(I-P_K)D_LP_K$ has kernel $K^\perp$, and its range
closure lies in $K^\perp$.
All these operators lie in the amplified equivariant algebra.
Polar decomposition there gives a partial isometry $U$ with
$U^*U=P_K$ and $UU^*\le I-P_K$.  One can obtain $U$ as the strong
limit of $T(T^*T+\varepsilon I)^{-1/2}$ as $\varepsilon\downarrow0$;
functional calculus and strong closure keep this limit in the algebra.
Therefore
\begin{equation}
 \tau(P_K)=\tau(U^*U)=\tau(UU^*)\le1-\tau(P_K),
 \qquad \tau(P_K)\le\frac12.
 \label{eq:dc-half-trace}
\end{equation}
This uses the range closure, not an assumption that the range is closed.
The expected root trace of $P_K$ is unchanged on forgetting the phases.

\emph{Return to finite matrices.}
Put $M_n=I-A_n$, $H_n=M_n^*M_n$, and $H=(I-A)^*(I-A)$.
Their normalized empirical spectral measures $\nu_n$ have common support
in $[0,c^2]$.  For each integer $j\ge0$, the moment
$(tN_n)^{-1}\operatorname{tr}H_n^j$ is the expected normalized root-block
trace of a bounded continuous local function; radius $2j$ suffices.
Local convergence gives moment convergence to the averaged root spectral
measure $\nu$ of $H$.  Polynomial approximation on $[0,c^2]$ gives weak
convergence.  The closed-set inequality for the singleton $\{0\}$ is
\begin{equation}
 \limsup_n\frac{\dim\ker M_n}{tN_n}
 =\limsup_n\nu_n(\{0\})
 \le\nu(\{0\})=\tau(P_K)\le\frac12.
 \label{eq:dc-nullity-semicontinuity}
\end{equation}
Here $\ker(M^*M)=\ker M$ holds for finite matrices and bounded operators.
This proves~\eqref{eq:dc-complex-rank}.  Real matrices have the same rank
after complexification.
\end{informalproof}

\begin{corollary}[Uniform finite form]
\notformalized
\label{cor:dc-uniform}
For fixed $\Delta,t,B$ and every $\varepsilon>0$ there exist an integer
$r\ge1$, a number $\delta>0$, and $N_0$ such that
\[
 b_E(r)\le\delta N,\quad N\ge N_0
 \quad\Longrightarrow\quad
 \operatorname{rank}_{\mathbb C}(I-A)\ge(1/2-\varepsilon)tN
\]
for every relation and matrix with the bounds of \cref{thm:dc-complex}.
\end{corollary}
\begin{informalproof}
Otherwise choose a failure with $r=j$, $\delta=1/j$, and $N\ge j$
for each positive integer $j$.  These failures have paired-distance separation,
contradicting \cref{thm:dc-complex}.
\end{informalproof}

\begin{corollary}[Complex linear codes on actual cone relations]
\notformalized
\label{cor:dc-native}
Fix the tape count, linear-time constant, full cone depth, block width~$t$,
and a uniform bound on the recovery coefficients.  Every sequence of
complex linear storage spaces on the corresponding native cone relations,
admitting linear recovery maps with the stated block width and coefficient
bound, has dimension at most $(1/2+o(1))tL$ as $d\to\infty$.
\end{corollary}
\begin{informalproof}
Remove the $O(L/\log d)$ self-demand coordinates supplied by
\cref{cor:nb-self-demands}.  For any fixed $0<\delta<1$, also remove
coordinates whose left degree exceeds a fixed $M\ge B_D/\delta$, where
$B_D=\sum_{h=0}^D(k+1)^h$ is the cone bound in
\eqref{eq:nb-cone-bound}, using
\cref{cor:nb-trim}.  Write $H$ for their union and $Q$ for the
retained set.  Then $|H|\le\delta L+o(L)$, and the induced relation is
loop-free, has bounded incidence degree, and satisfies the paired-distance
condition by \cref{thm:nb-main}.

If $\mathcal U$ is the original linear space, its subspace with zero
coordinates on $H$ has codimension at most $t|H|$.  On this subspace the
original bounded linear recovery maps restrict to maps on $Q$ with the
same coefficient bound and block width.  \cref{thm:dc-complex} bounds
its dimension by $(1/2+o(1))t|Q|$.  Therefore
\[
 \dim\mathcal U\le t|H|+(1/2+o(1))t|Q|
 \le(1/2+\delta/2+o(1))tL.
\]
Take the size limit first and then let $\delta$ tend to zero.  This is a
complex-linear statement about the cone relations; it does not make the
machine's binary genie a complex linear map.
\end{informalproof}

\begin{remark}[Scope of the graph extension]
\label{rem:dc-comparison}
For a symmetric recovery relation, $E$ is its bipartite double cover.
A path between matched ports projects to an odd closed walk, and an odd
cycle gives such a path at each of its vertices.  With bounded degree,
paired-distance separation is consequently equivalent to having only $o(N_n)$
vertices on odd cycles of every fixed bounded length.
For directed relations the hypotheses differ: in disjoint directed
seven-cycles with one recovery neighbour at each vertex, $E$ is a union
of isolated edges and every matched distance is infinite, whereas the
underlying coordinate graphs have odd cycles at positive density.
Thus replacing directed recovery by its undirected closure loses information.
\end{remark}

\subsection{What the complex bound says over finite fields}

\begin{corollary}[Necessary torsion in bounded lifts]
\notformalized
\label{cor:dc-torsion}
Fix a prime $p$, integers $\Delta,t\ge1$, and $\varepsilon>0$.
Let $(E_n)$ satisfy the graph hypotheses of \cref{thm:dc-complex}.
Suppose matrices $M_n$ over $\mathbb F_p$ have diagonal blocks $I_t$,
off-diagonal support in $E_n$, and
$\dim\ker M_n\ge(1/2+\varepsilon)tN_n$.
For any integer lifts $\widetilde M_n$ with the same support restrictions,
literal diagonal blocks $I_t$, and uniformly bounded entries, at least
$(\varepsilon-o(1))tN_n$ of their \emph{nonzero} Smith invariant factors
are divisible by $p$.
\end{corollary}
\begin{informalproof}
Set $\rho_n=\operatorname{rank}_{\mathbb Q}\widetilde M_n$.
\cref{thm:dc-complex} gives $\rho_n\ge(1/2-o(1))tN_n$.
The unimodular transformations in Smith normal form remain invertible
modulo $p$.  Thus
\[
 \#\{\hbox{nonzero invariant factors divisible by }p\}
 =\rho_n-\operatorname{rank}_{\mathbb F_p}M_n
 \ge(\varepsilon-o(1))tN_n.
\]
Equivalently, the subgroup annihilated by $p$ in the integer cokernel
has at least this dimension over $\mathbb F_p$.  Zero invariant factors
give free abelian summands and are not counted as torsion.
\end{informalproof}

This conclusion permits such torsion; it is not a finite-field half-rate
theorem.  A fixed extension field can be represented over its prime field
by enlarging the fixed block width.  The analytic argument does not cover
growing block widths, unbounded complex coefficients, or nonlinear codes
on nonhyperfinite relations.  \cref{sec:mc-section} gives a separate
finite-field certificate on one existential expanding family, uniform
in the storage-code block width.  Neither result identifies all native
cone relations with that family or settles the unrestricted coding problem.

\section{A native nonhyperfinite family with a half-rate bound}
\label{sec:ne}

The directed recovery formulation does not require a decomposition into
small components.  We demonstrate this with one fixed native computation.
Its full cone incidence graphs contain robust expander cores, yet every
storage code on every induced coordinate set has logarithmic cardinality
at most half its coordinate count plus $o(L)$.  This gives rate at most
one half plus $o(1)$ whenever the retained set has size $\Omega(L)$.
The computation is an invertible shear,
not a transposition algorithm.  The point of the example is the simultaneous
graph-theoretic and coding conclusion for its actual execution.

Throughout this section let
\[
 d=2m,\qquad b=\max\{1,\lceil\log_2d\rceil\},\qquad
 W=db,\qquad L=dW=d^2b.
\]
A raw coordinate is $p=(i,j,s)\in V=[d]\times[d]\times[b]$, where
$[u]=\{0,\ldots,u-1\}$.  Its physical address is $(id+j)b+s$, and its
transpose demand is the physical output address of
$\Pi(p)=(j,i,s)$.  The paired incidence graph has distinct ports $L_p,R_p$;
an edge $L_pR_q$ means that raw input $p$ occurs in the full backward cone
of the append at $\Pi(q)$.  We write this graph as $E_\Gamma$ when cones
are stopped by a whole-step genie $\Gamma$.  Deleting a coordinate always
deletes both of its ports.  In particular it is different from deleting
chosen dependencies.

\subsection{Two imported interfaces and a paid executor}

We use the following expander input from the native shear construction of
\cite[section on expander shears and local recovery]{BrendeLocalCoding}.
The graph construction is imported; the schedule refinements and the
subsequent cone analysis will be given here.

\begin{lemma}[Ordered expander input]\label{lem:ne-expander-import}
\cited{The expander-shear construction of~\cite{BrendeLocalCoding}. The uniform list-generation refinement is justified below by a written argument; this full interface is not claimed to be Lean-verified.}
There are fixed positive constants $J,\Delta,\eta$, with $J,\Delta$
integers and $\eta\le1$, and simple bipartite graphs $H_m$ on
\[
 A_0,\ldots,A_{m-1},\quad B_0,\ldots,B_{m-1},
\]
such that their maximum degree is at most $\Delta$, their edges partition
into $J$ partial increasing matchings, and
\[
 |N_{H_m}(S)\setminus S|\ge\eta|S|
 \quad\text{whenever }|S|\le m.
\]
The matching endpoint lists can be generated uniformly in
$O(m\operatorname{polylog}m)$ time using a fixed number of tapes and
polylogarithmic working space besides the monotonically emitted lists.
All constants and rational coefficients are fixed before $m$ varies.
\end{lemma}

Here a partial increasing matching has its two endpoint sequences strictly
increasing.  The computational part of the imported interface can be seen
directly from the rational overlap construction.  Fix its finite family
of increasing rational maps on interval domains.  Their derivatives have
fixed positive lower and upper bounds.  An edge records a positive-length
overlap of the image of one mesh cell with a second mesh cell; exact
rational comparisons decide boundary coincidences.  There are constantly
many candidate neighbors in either direction.  Split candidates by map
name and by the residues of the two endpoint indices modulo a fixed
integer larger than these consecutive-neighbor bounds.  Each resulting
class is an increasing matching.  Symmetrization uses further fixed
classes, and retaining a duplicate only in its first class preserves the
simple graph.  Evaluation, division, comparison, and duplicate tests use
polynomially many bit operations and cells in $\log m$ per candidate.
Enumerating sources and filtering by class emits each page in increasing
order.  Thus this use of the expander input requires neither a table
depending on $m$ nor a unit-cost operation on a growing word.

The depth interface is a machine-independent consequence of the previous
paper.  We state precisely the portion needed below.

\begin{lemma}[Cheap whole-step depth for home-based walks]
\cited{Theorem B and the proof of Corollary B of~\cite{BrendeCurved}. Both the home-based cut theorem and the whole-step conversion used here are written proofs, not Lean-verified results.}
\label{lem:ne-depth-import}
Fix an oblivious native machine with $k$ work tapes, whole-step value width
at most $\beta$ bits, and running time at most $c(L+1)$.  Suppose every
work-head walk is home-based: every greedy weakly monotone leg begins or
ends within distance one of its fixed home cell.  For every fixed
$\varepsilon>0$, and all sufficiently large admissible $L$, there is a
whole-step genie satisfying $\beta|\Gamma|\le\varepsilon L$ whose full
depth, in the sense of \cref{def:nb-cones}, is at most a constant $D$.
The constant may
depend on the machine and $\varepsilon$, but not on $L$.
\end{lemma}

This is the cut theorem and long-edge conversion of
\cite[Theorem B and the proof of Corollary B]{BrendeCurved}, before any
assumption of transposition correctness.  A cut edge $t\to s$ is blocked
by revealing $s$; a periodic set of additional revealed steps bounds the
runs of state edges and pauses between long hops.  The imported statement
therefore controls all edges of the whole-step DAG, not just long hops.
We will use it only for existence.  Our conclusions about a specified
depth-$D$ genie quantify over every such genie.

\begin{proposition}[A fixed native shear executor]\label{prop:ne-native-shear}
\notformalized
There is a fixed deterministic oblivious native machine $M_0$ which, at
all sufficiently large even canonical dimensions, computes
\[
 \Phi(A,B)=(A,B+L_mA)
\]
in $O(L)$ transitions, where $A,B$ are the first and last $m$ raw rows and
$L_m$ is the binary adjacency matrix of $H_m$.  It uses read-only raw
input, blank initial work tapes, and append-only raw output, and all work
heads are home-based.  There is also such a machine $M_1$ which appends
all $A$ rows only after completing the shear pages, followed by all $B$
rows.  Both machines have fixed alphabet, state set, and tape count.
\end{proposition}

\begin{informalproof}
We give the execution details because later arguments use its visits,
including visits at which the scanned symbols are semantically irrelevant.
Two separate work tapes will hold banks of length $n=mW$.  They begin
blank; each receives a persistent origin flag, row-start flags during
loading, and an end marker when its loading finishes.  A unary row clock
has length $W$ and its own origin and end marker.  All flags are tracks
of a fixed finite alphabet.  Both data heads remain at their origins
throughout setup, so setup does not visit or initialize the full banks.

First count the raw input length with a little-endian binary counter whose
head returns to its origin after each carry.  The sum of carry lengths is
$O(L)$, including return traversals.  The raw input is then rewound.
Canonical shape recovery uses polynomially many bit operations in
$\log L$: try the possible entry widths, compute an integer square root,
and check the defining equality.  Invalid lengths, odd dimensions, and a
fixed finite set of small cases use an identity fallback.  Only the even
canonical subfamily is used in this section.

All polylogarithmic-space arithmetic can be made physically home-based
at polynomially larger polylogarithmic cost.  To make this explicit,
bootstrap a unary scratch ruler of length $S=(\lceil\log_2(L+2)\rceil+c)^h$
by repeated complete sweeps.  Here the fixed exponent $h$ exceeds the
space exponent of the routines to be simulated.  A fixed number of tracks
in this window represent virtual tapes, symbols, and virtual head flags.
A complete outward-and-return scan reads the flagged symbols; further
complete scans implement a virtual transition and relocate flags.  Right
moves can be flagged on an outward scan and left moves on a return scan,
using old and new flag tracks.  Each virtual transition costs $O(S)$,
and every physical turn occurs at the origin or at the end of the ruler.
The length counter itself is implemented by the direct carry routine,
so this simulation does not multiply its linear cost by a polylogarithm.

Generate the $J$ endpoint lists of \cref{lem:ne-expander-import} as
monotone output streams.  Copy an individual record to the origin buffer
before parsing it.  The long list heads make monotone traversals and full
rewinds, rather than reversing locally for arithmetic.  Their total length
is $O(m\log m)$ and the compiler still costs
$O(m\operatorname{polylog}L)$.  Construct the unary row clock by reading
the binary digits of $W$ from most to least significant.  At each digit,
double the current unary string onto an alternate tape, append the digit,
erase and rewind the old string, and rewind the new one.  The geometric
sum of the string lengths is $O(W)$.  All these work heads are home-based.

Load the $A$ bank and then the $B$ bank in raw row order.  The row clock
identifies the first bit of each row, where the load writes the row-start
flag; completion of a bank writes its end marker at $n$.  Thus these
flags require no preliminary full-bank sweep.  An active bank head and
the row-clock head advance one cell per bit.  At a row boundary
the bank head waits while the clock returns to its origin.  The machine
$M_0$ also appends each $A$ bit at its load transition; $M_1$ omits these
early appends.  After loading, rewind both banks completely.

For each increasing matching page, read its next endpoint pair $(i,j)$,
skip whole rows forward until the two indicated rows are reached, and
scan the pair in alignment.  In a single transition at offset $s$, read
$A_i[s]$ and $B_j[s]$, write their XOR to $B_j[s]$, and advance both heads
and the row clock.  A skip likewise uses one transition per bit and
preserves its payload.  Clock rewinds and controller operations take place
with the bank heads parked at row starts or at the bank end.  Increasing
endpoints make a single forward traversal sufficient.  After the last
matched pair, scan the unused suffixes to the end and rewind both banks
in lockstep, one cell per transition.  Empty pages use the same complete
traversal convention.  Each page costs $O(L)$.

All work cells scanned on every transition are rewritten, including cells
under parked heads.  Preserving a symbol is a write, and no part of the
construction assumes otherwise.  Induction over the pages shows that the
$B$ bank contains $B+L_mA$: every simple edge is executed exactly once,
so duplicate XORs do not cancel a desired contribution.  The $A$ bank
remains unchanged.  Finally $M_0$ appends the $B$ bank.  For $M_1$, after
the last full rewind, scan and append $A$ and then $B$.  This has the
specified delayed output order.

The complete cost is
\[
 O\bigl(L+W+JL+m\operatorname{polylog}L+\operatorname{polylog}L\bigr)
 =O(L),
\]
since $L=4m^2b$ and all logarithmic powers are fixed.  Public parameters,
flags, and endpoint records determine every move and append time; payload
values determine only written or appended bits.  The schedule is therefore
oblivious.  Bank and metadata heads make complete outward traversals and
returns, possibly pausing along a traversal.  Clock and compiled controller
heads also make origin excursions.  These are home-based physical walks.
The number of pauses need not be bounded, and we make no assertion of a
bounded number of affine pieces within a leg.
\end{informalproof}

\subsection{Literal expander paths and robust separators}

The native DAG has a state edge $t\to t-1$ and, for each work tape, an
edge from $t$ to the previous visit to its scanned cell, if one exists.
Its in- and out-degrees are at most $k+1$ after parallel edges are merged.
Raw input bits attach to their read steps; the read-only input tape does
not contribute previous-visit edges.  A genie-free path may start at a
revealed vertex but has no subsequent revealed vertex.  When proving
survival below, we conservatively discard a path meeting the genie even
at its first vertex.

\begin{lemma}[A literal core with bounded congestion]\label{lem:ne-paths}
\notformalized
For either machine in \cref{prop:ne-native-shear}, every edge
$A_iB_j$ of $H_m$ and every lane $s\in[W]$ has a literal backward DAG
path from the append of $B_j[s]$ to a read of raw input $A_i[s]$.
The chosen paths have maximum vertex congestion at most $C_0=2\Delta$,
including endpoints.  For lanes $1\le s<W$, their lengths are bounded by
a constant depending only on $J$.  No length bound is needed for lane zero.
\end{lemma}

\begin{informalproof}
Let $t_{\mathrm{in}}$ be the load of $A_i[s]$, $t_{\mathrm{xor}}$ the update for the
chosen edge, and $t_{\mathrm{out}}$ the append of $B_j[s]$.  List every visit
to the physical $A_i[s]$ cell from $t_{\mathrm{in}}$ through $t_{\mathrm{xor}}$;
then list every visit to the $B_j[s]$ cell from $t_{\mathrm{xor}}$ through
$t_{\mathrm{out}}$.  Consecutive visits in each list are actual previous-visit
edges when read backward, and the lists share the XOR vertex.  Their
concatenation gives the claimed path to an actual raw read.  It includes
unchanged writes and controller waits.

For $s\ne0$, no bank head parks at this cell.  The $A$ cell has at most
$2J+2$ relevant visits and the $B$ cell at most $2J+3$; a bound of
$4J+8$ on path length is more than sufficient.  At lane zero, all parked
visits are included and can make the path long.  At any time a scanned
$A$ cell belongs to paths for at most $\Delta$ incident expander edges,
and the same holds for the scanned $B$ cell.  Their union has size at
most $2\Delta$.  The extra delayed $A$ scan in $M_1$ can add parked
visits on the $B$ origin, but does not change this congestion estimate.
\end{informalproof}

Thus $E_\Gamma$ contains $W$ vertex-disjoint copies of $H_m$ with at most
$C_0|\Gamma|$ missing core edges.  There are $2mW=L$ core ports, a subset
of the $2L$ incidence ports.  The claim concerns containment in the full
syntactic cone graph; unselected state and cell paths may supply further
edges.

\begin{proposition}[A robust separator estimate]\label{prop:ne-separator}
\notformalized
Let $U$ be any deleted coordinate set.  If a further port set $Z$ makes
every component of $E_\Gamma-U-Z$ have size at most $K\le m$, then
\[
 |Z|\ \ge\ \frac{\eta L-2C_0|\Gamma|}{\Delta+\eta}-2|U|.
\]
This holds for every genie, without a depth hypothesis.
\end{proposition}

\begin{informalproof}
Discard the at most $f=C_0|\Gamma|$ core edges whose paths meet $\Gamma$.
At most $z\le2|U|+|Z|$ core ports have been deleted.  Intersect each
remaining component with each lane core.  These intersections have size
at most $m$, so expansion applies to each.  Summing their edge-boundary
lower bounds counts at least $\eta(L-z)$ crossing edges with multiplicity.
An edge joining two surviving intersections must be one of the $f$ lost
core edges, and is counted at most twice.  Edges meeting the deleted core
ports contribute at most $\Delta z$.  Hence
\[
 \eta(L-z)\le2f+\Delta z.
\]
Rearrange and use $z\le2|U|+|Z|$.  Additional cone edges can merge
components, but cannot invalidate this lower bound.
\end{informalproof}

\subsection{The full syntactic cones retain a lane constraint}

The preceding core alone says nothing about code rate.  For that purpose
we must control every cone edge, including those introduced by parked
heads and state transitions.  Fix either $M_0$ or $M_1$, let $k$ be its
number of work tapes, and write
\[
 B_D=\sum_{h=0}^D(k+1)^h.
\]

\begin{lemma}[A sparse interface for full cones]\label{lem:ne-lane}
\notformalized
There are fixed constants $C_s,h_s$ and a set $X$ of actual time vertices,
depending only on the length-specific schedule, such that
\[
 |X|\le S_d:=C_s\bigl[W+d(1+\log_2(L+2))^{h_s}\bigr].
\]
For every full-depth-$D$ genie, every output outside the $D$-step forward
neighborhood of $X$, and every non-row-start raw input in its cone, the
two physical within-row bit offsets have cyclic distance at most $D$
in $\mathbb Z/W\mathbb Z$.
\end{lemma}

\begin{informalproof}
The entire length count, shape calculation, metadata generation, and clock
construction form the setup.  This can take $\Theta(L)$ transitions, so
it cannot be discarded as a small set of times.  Its distinct work cells,
however, number only $O(W+d\operatorname{polylog}L)$.  Mark the first
post-setup state vertex and the first post-setup visit to every cell touched
during setup.  Any backward state or cell edge entering setup begins at
one of these marked vertices.

Post-setup fast phases are a forward row load, skip, or XOR scan; a clock
return; a simultaneous whole-bank rewind; and a final forward output
scan.  Mark all other controller, parser, arithmetic, and metadata ticks.
Their total is $O(d\operatorname{polylog}L)$.  Mark also the first and last
ticks of each fast phase, moving bank visits to offset zero modulo $W$,
and moving clock visits to its endpoints.  There are $O(d)$ such marks,
with the fixed number of pages absorbed in the constant.  A stationary
head at offset zero is not marked at every tick.

Every unmarked post-setup vertex has a lane: the position of its moving
clock modulo $W$, if the clock moves, and otherwise that of its moving
bank.  Simultaneously moving banks have the same lane.  A state edge
between unmarked vertices changes lane by at most one.  For a cell edge,
a stationary head has its immediately preceding time as predecessor once
the first fast tick has been excluded.  A moving bank head at a nonzero
offset has no parked visits at that cell; its preceding post-setup moving
visit has the same lane.  The same is true of a moving interior clock
cell, in a forward row scan or return.  These cases cover all work tapes,
including tapes whose fields the current operation ignores.  Consequently
every edge between good vertices changes lane by at most one cyclically.

The $D$-step forward neighborhood of $X$ contains at most $B_D|X|$ time
vertices, since forward degree is at most $k+1$.  An output outside it
has no marked vertex in a depth-$D$ free cone.  Its cone also cannot enter
setup, by the first-touch marks.  During loading, a non-row-start raw bit
is read only at a load vertex with that same lane.  During row-boundary
pauses the raw head holds the next row-start bit or is outside the payload;
after loading it is outside the payload.  Excluding the $d$ raw row starts
therefore accounts for every remaining input attachment.  Summing lane
changes along a free path proves the result.  The output itself is included
in the forward-neighborhood exclusion, which covers the convention allowing
a free path to start at a revealed vertex.
\end{informalproof}

The following elementary code argument will be useful for both examples.

\begin{lemma}[An order of port colors]\label{lem:ne-color-injection}
\formal{ColorOrder.lean}{color_order_half_bound}
Let a paired incidence graph on $N$ coordinates admit colors
$\alpha(v)$ at $L_v$ and $\zeta(v)$ at $R_v$ such that every edge is
monochromatic and $\alpha(v)\ne\zeta(v)$ for every $v$.  Every storage
code over an alphabet of cardinality $a$ has size at most
$a^{\lfloor N/2\rfloor}$.
\end{lemma}

\begin{outline}
Order the finitely many used colors.  Record the coordinates $v$ for which
$\alpha(v)$ precedes $\zeta(v)$.  Process left-color classes in increasing
order.  If a coordinate was not recorded, then its right color precedes
its left color.  Every coordinate in its recovery neighborhood has that
earlier left color, by monochromaticity, and so has already been recovered.
This proves that projection onto the recorded coordinates is injective on
the code; no extension of an arbitrary partial assignment is needed.
Compare the color order with its reverse. Each coordinate is recorded in
exactly one of them, so one record has at most $\lfloor N/2\rfloor$
coordinates, which proves the claim.
\end{outline}

\begin{theorem}[Half rate on every induced full-cone relation]
\notformalized
\label{thm:ne-full-cone-half}
Fix $D\ge0$.  For either native machine in
\cref{prop:ne-native-shear}, for every whole-step genie of full
depth at most $D$, every $Q\subseteq V$, and every storage code
$\mathcal C\subseteq\mathcal A^Q$ for the induced actual cone relation,
where $a=|\mathcal A|\ge2$ and $N=|Q|$, there is $u\le\min(N,e_D(d))$
such that
\[
 |\mathcal C|\le a^{\lfloor(N+u)/2\rfloor}.
\]
One permissible error bound, uniform in the genie, alphabet, and $Q$, is
\[
 \begin{split}
 g&=\lceil\sqrt d\rceil,\qquad r=\lceil d/g\rceil,\\
 e_D(d)&=B_DS_d+d+2d(2D+1)r+bdg=o_D(L).
 \end{split}
\]
There is no bound on the cost of the genie in this theorem.
\end{theorem}

\begin{informalproof}
Partition the column indices into $r$ consecutive buckets of at most $g$
columns each, and view their bit lanes as consecutive arcs of
$\mathbb Z/W\mathbb Z$.  Guard all lanes within cyclic distance $D$ of
a bucket boundary.  There are at most $(2D+1)r$ guarded bit lanes.
Delete from $V$ the following coordinate sets: demand labels of the bad
physical outputs in \cref{lem:ne-lane}; the $d$ raw row starts; coordinates
whose input or transpose-output lane is guarded; and coordinates $(i,j,s)$
whose row and column belong to the same bucket.  The first set costs at
most $B_DS_d$, the third at most $2d(2D+1)r$, and the fourth at most $bdg$.
The total deletion is at most $e_D(d)$.

On the surviving ports give $L_{(i,j,s)}$ the bucket of $j$, and give
$R_{(i,j,s)}$ the bucket of $i$.  A surviving incidence edge has equal
colors by \cref{lem:ne-lane}: a path of lane length at most $D$ cannot
cross a bucket boundary while both endpoint lanes avoid its guard.
The own-port colors differ by the fourth deletion.  The graph therefore
satisfies \cref{lem:ne-color-injection}.

Intersect the deleted set with $Q$, obtaining $u$ coordinates.  Condition
their values in the original code.  Each nonempty fibre is a storage code
on the induced surviving relation: the removed recovery inputs have fixed
values in that fibre.  \cref{lem:ne-color-injection} bounds its size by
$a^{\lfloor(N-u)/2\rfloor}$.  There are at most $a^u$ fibres, giving
$a^{u+\lfloor(N-u)/2\rfloor}=a^{\lfloor(N+u)/2\rfloor}$.
For fixed $D$, each displayed error divided by $L=d^2b$ tends to zero.
All exclusions were chosen from the schedule and the full relation; the
argument is simultaneous for all alphabets, recovery functions, and codes.
\end{informalproof}

\begin{corollary}[Nonhyperfinite bounded-degree cones with half rate]
\notformalized
\label{cor:ne-combined}
For one fixed native home-based linear-time machine there are constants
$\varepsilon,\alpha>0$ and fixed integers $D,M$ with the following
property.  At every sufficiently large even canonical dimension the class
of genies satisfying $\beta|\Gamma|\le\varepsilon L$ and full depth at
most $D$ is nonempty.  For every such genie, one can delete coordinates
to obtain an induced paired graph on $N=\Theta(L)$ coordinates, of maximum
degree at most $M$, with no own-port edge, such that:
\begin{enumerate}
\item for every fixed $K$ and all sufficiently large $d$, every port set
whose deletion leaves components of size at most $K$ has size at least
$\alpha L/2$;
\item every storage code on every further induced coordinate set has
logarithmic cardinality at most half its coordinate count plus $o_D(L)$,
uniformly over all finite alphabets.
\end{enumerate}
\end{corollary}

\begin{informalproof}
Choose $\gamma>0$ with $\gamma\le\eta/(8C_0)$, and then
$0<\varepsilon\le\min(1,\beta\gamma)$.  Set
$\alpha=\eta/[2(\Delta+\eta)]$.  \cref{prop:ne-separator}
gives a separator lower bound at least $\alpha L$ before coordinate
trimming for every genie of the stated price.  \cref{lem:ne-depth-import}
supplies a fixed $D$ and makes the permitted genie class nonempty.

Each right incidence degree is at most $B_D$, so there are at most
$B_D L$ edges.  Choose a constant degree threshold $M_0\ge8B_D/\alpha$.
Delete every coordinate whose left degree exceeds $M_0$; their number is
at most $\alpha L/8$.  Also delete coordinates with an own-port edge.
The one-stage transpose-matching estimate of \cref{cor:nb-self-demands}
bounds these by $O(L/\log d)$, uniformly over the genies in question;
eventually this is at most $\alpha L/8$.  This use of the matching estimate
does not require transpose correctness or an initial left-degree bound.
The remaining graph has maximum degree at most $M=\max(M_0,B_D)$ and
$N=\Theta(L)$.  The separator loss is at most twice the number of deleted
coordinates, giving $\alpha L/2$.  For any fixed $K$, eventually $K\le m$,
so \cref{prop:ne-separator} applies.  The code conclusion is
\cref{thm:ne-full-cone-half}, whose quantifiers already include every
induced coordinate set.
\end{informalproof}

\begin{remark}\label{rem:ne-scope}
The conclusion rules out a uniform small-component decomposition for this
family at arbitrarily small deletion cost.  It does not exhibit a high-rate
code.  In fact its simultaneous half-rate bound shows why a directed
coding theorem can remain possible after a separator approach fails.
Neither the lane lemma nor its coloring is claimed for arbitrary native
schedules.  The next construction changes the paid physical layout and
destroys almost-everywhere separation of paired ports by component colors.
\end{remark}

\section{Connected demand pairs and a depth-independent coding deficit}
\label{sec:cd}

The coloring in the preceding section separates the two ports of almost
every coordinate.  We now give a paid native modification for which no
such separation by connected components is possible after a sublinear
coordinate deletion.  Nevertheless its full cone relations have a
positive all-code deficit whose coefficient is independent of the depth.
The proof uses a finite row-phase description and an independent rectangle,
and does not recover a half-rate bound.

We retain $d=2m$, $b=\max(1,\lceil\log_2d\rceil)$, $W=db$, $L=d^2b$,
the graphs $H_m$, and the physical transpose-demand labels of
\cref{sec:ne}.  Indices of $B$ rows in this section are global:
$m\le j<d$.  Define a row-dependent rotation by
\[
 R:(i,c,s)\longmapsto(i,c+i\pmod d,s).
\]
The new machine computes the invertible map $R^{-1}\Phi R$.
Throughout this section, $k$ and $\beta$ denote the work-tape count and a
whole-step value-width bound for this rotated machine $M_R$, including
its additional buffers and controller tapes.  They need not equal the
corresponding bounds for the unwrapped machines.  We use
$B_D=\sum_{h=0}^{D}(k+1)^h$ with these enlarged resources.

\subsection{A robust count of connected pairs}

\begin{lemma}[Prescribed pairs in damaged expanders]
\formal{DamagedExpanders.lean}{DamagedExpanders.damaged_expander_copies}
\label{lem:cd-connected-pairs}
Let $H_1,\ldots,H_W$ be disjoint graphs on $2m$ vertices, each of maximum
degree at most $\Delta$ and with outside vertex expansion at least
$\eta|S|$ for every $|S|\le m$, where $\eta>0$.  Prescribe $n_*$ pairs of distinct
vertices, with all $2n_*$ endpoints distinct and each pair contained in
one graph.  After deletion of $f$ edges and $z$ vertices, at least
\[
 n_*-\frac f\eta-\left(1+\frac\Delta\eta\right)z
\]
prescribed pairs survive with their endpoints in the same connected
component.  The bound holds for every finite $m\ge1$ and arbitrary
distributions of the deletions among the graphs.
\end{lemma}

\begin{outline}
Vertex expansion implies an edge-boundary bound $|\partial_E S|\ge\eta|S|$.
Consider one graph, and write $f_0,z_0$ for its deletion counts.  If the
surviving graph has a component $G$ of more than $m$ vertices, let $S$ be
the union of its other surviving components.  Then $|S|<m$.  Every original
edge from $S$ to $G$ was deleted, and at most $\Delta z_0$ edges go from
$S$ to deleted vertices.  Thus
\[
 \eta|S|\le f_0+\Delta z_0.
\]
Every prescribed pair wholly inside $G$ is connected.  All other pairs
can be charged to distinct endpoints in $S$ or the deleted set, since the
prescribed endpoints are disjoint.  Their number is at most
$f_0/\eta+(1+\Delta/\eta)z_0$.

If every surviving component has size at most $m$, sum the edge-boundary
inequality over those components.  Edges between surviving components
were deleted and are counted at most twice; edges incident to a deleted
vertex contribute at most $\Delta z_0$.  Hence
\[
 \eta(2m-z_0)\le2f_0+\Delta z_0,
 \qquad
 m\le\frac{f_0}{\eta}+\frac{\Delta+\eta}{2\eta}z_0.
\]
There were at most $m$ prescribed pairs in this graph.  Declaring all of
them lost therefore satisfies the same asserted bound.  This includes an
empty surviving graph.  Summing over the copies proves the lemma.
\end{outline}

\subsection{A physical rotation wrapper}

\begin{proposition}[The rotated native shear]\label{prop:cd-rotated-native}
\notformalized
There is one fixed deterministic oblivious native machine $M_R$ which
computes $R^{-1}\Phi R$ in $O(L)$ transitions at all sufficiently large
even canonical dimensions.  All work heads are home-based.  For each
edge $A_iB_j$ of $H_m$ and each internal lane $(g,s)\in[d]\times[b]$
there is a literal backward path from the final physical append
\[
 (j,g-j\pmod d,s)
 \quad\hbox{to a raw read at}\quad
 (i,g-i\pmod d,s).
\]
These paths have endpoint-inclusive whole-step vertex congestion at most
$C_1=4\Delta$.  Their lengths need not be bounded.
\end{proposition}

The complete native verification is in \cref{sec:details-cd-rotated-native}.


\begin{theorem}[Many robustly connected demand pairs]
\notformalized
\label{thm:cd-robust-demands}
For $M_R$, every genie $\Gamma$ and every deleted coordinate set $U$
leave at least
\begin{equation}\label{eq:cd-pair-bound}
 \frac L4-\frac{C_1}{\eta}|\Gamma|
             -2\left(1+\frac\Delta\eta\right)|U|
\end{equation}
surviving coordinates whose two ports are connected in the induced full
cone graph.  There are constants $\varepsilon>0$ and fixed integers
$D,M$ for which the permitted cheap depth-$D$ genies are nonempty and,
for every such genie, the usual degree and own-port-edge coordinate
trims yield a bounded-degree graph on $N=\Theta(L)$ coordinates with at
least $L/8$ connected demand pairs.  After any additional $o(L)$
coordinate deletion at least $L/16$ such pairs remain for all sufficiently
large $d$.
\end{theorem}

\begin{informalproof}
For each internal lane $(g,s)$, the literal paths give a copy of $H_m$
with core ports
\[
 A_i\longleftrightarrow L_{(i,g-i,s)},\qquad
 B_j\longleftrightarrow R_{(g-j,j,s)},
\]
where columns are reduced modulo $d$.  These $W$ cores have $L$ distinct
ports.  For a coordinate $p=(i,c,s)$ with $i<m$ and $c\ge m$, its two
ports lie in the same core lane $g=i+c\pmod d$.  Prescribe these
$m^2b=L/4$ disjoint pairs.  Their distribution among lanes need not be
uniform.

Discard every selected core path meeting $\Gamma$, including endpoints.
At most $C_1|\Gamma|$ edges are discarded.  Every remaining path witnesses
an actual full-cone incidence.  Deleting $U$ removes at most $2|U|$ core
ports.  \cref{lem:cd-connected-pairs} proves~\eqref{eq:cd-pair-bound},
and extra full-cone edges can only join components.

Put $A_*=C_1/\eta$ and $B_*=2(1+\Delta/\eta)$.  Choose
$0<\varepsilon\le\min(1,\beta/(16A_*))$.  For every genie of price
$\beta|\Gamma|\le\varepsilon L$, the first loss in
\eqref{eq:cd-pair-bound} is at most $L/16$.  \cref{lem:ne-depth-import}
supplies a fixed full depth $D$ and nonemptiness.  Fix this $D$ and let
$B_D=\sum_{h=0}^D(k+1)^h$.  Choose $M_0\ge32B_*B_D$.
There are at most $B_D L$ incidence edges, so deleting coordinates with
left degree above $M_0$ costs at most $L/(32B_*)$.  The self-demand
estimate of \cref{cor:nb-self-demands} bounds the own-port-edge deletion by
$O(L/\log d)$, eventually also at most $L/(32B_*)$.  The remaining
maximum degree is at most $M=\max(M_0,B_D)$ and its coordinate count is
at least $L-L/(16B_*)$.  Formula~\eqref{eq:cd-pair-bound} now leaves
at least $L/4-L/16-L/32-L/32=L/8$ pairs.  An additional set $U'$ costs
at most $B_*|U'|$ in that formula, proving the last assertion.

The choices have the order machine, $\varepsilon$, $D$, $M_0$, and
finally sufficiently large $d$.  None of $D,M_0$ grows with dimension.
All conclusions hold for every genie satisfying the fixed price and
depth bounds, rather than for a selected favorable genie.
\end{informalproof}

An edge-monochromatic port coloring is constant on each connected
component.  Thus \cref{thm:cd-robust-demands} precludes separating
all surviving own-port pairs by such a coloring after $o(L)$ further
deletion, even after paying the fixed degree and self-demand trims.
This does not contradict the directed incidence bridge: for every fixed
radius, that bridge bounds the number of short paired paths by
$O(L/\log d)$.  Almost all of the connected pairs above consequently
have distance exceeding any prescribed fixed radius.  Connectivity and
bounded-distance recovery are different assertions.

\subsection{Every full cone has a fixed row-phase description}

We next analyze all syntactic dependencies of the same machine, including
the auxiliary tapes.  We do not infer an incidence upper bound from the
selected expander paths or from the support of the shear matrix.

\begin{lemma}[A fixed number of row words]\label{lem:cd-row-words}
\notformalized
For the complete trace of $M_R$ there are fixed integers $K,\ell_{\max}$ and fixed
constants $C,h$ with the following property.  For every fixed $D\ge0$
and every full-depth-$D$ whole-step genie, there is a coordinate exception
set $Z$ with
\[
 |Z|\le z_D(d):=B_D C\bigl[W+d(1+\log_2(L+2))^h\bigr]+d=o_D(L)
\]
such that every actual cone incidence $L_pR_q$ with $p,q\notin Z$ has
its physical output row $\operatorname{outrow}(q)$ and raw input row
$\operatorname{rawrow}(p)$ related by one of at most $K$ row words.
Each word has at most $\ell_{\max}$ strictly increasing partial injections and
bands of radius $r=D+2$.  Its direction is from physical output row to
raw input row.  One may take
\[
 M_{\mathrm{ph}}=2J+3,\qquad K=3^{M_{\mathrm{ph}}}+1,
 \qquad \ell_{\max}=2(M_{\mathrm{ph}}-1).
\]
The maps can depend on the dimension's actual scan events; their number
and the word lengths do not depend on $d,D$, or the genie.
\end{lemma}

The complete native verification is in \cref{sec:details-cd-row-words}.


\subsection{Applying the row-word deficit}

\begin{theorem}[A positive deficit independent of depth]
\partformal{RowWordDeficit.lean}{row_word_cover_independent,row_word_cover_code_log_bound,row_word_cover_asymptotic_code_gap}{The native row-word extraction and its uniform sublinear exceptional-count estimate.}
\label{thm:cd-uniform-gap}
There is a constant $\rho>0$, depending only on the fixed machine $M_R$,
with the following property.  For every fixed $D$, every whole-step genie
of full depth at most $D$, every induced coordinate set $Q\subseteq V$,
and every nonempty storage code $\mathcal C\subseteq\mathcal A^Q$ for
the actual full-cone relation, put $N=|Q|$, $a=|\mathcal A|\ge2$.
Then
\[
 \log_a|\mathcal C|
 \le N-\max\{0,L_D(d)-|V\setminus Q|\},
\]
where, for all $d\ge4K$, one may take
\[
 \begin{split}
 t_D&=K(4K-1)\bigl[(2D+5)^{\ell_{\max}+1}-1\bigr],\\
 L_D(d)&=\max\left\{0,
 b(2K+1)\left\lfloor\frac d{4K}\right\rfloor^2
                    -bd\min(d,t_D)-z_D(d)\right\},
 \end{split}
\]
with $K,\ell_{\max},z_D$ from \cref{lem:cd-row-words}.  In particular one can
choose $\rho=1/(8K)$, independently of $D$, and obtain
\[
 \log_a|\mathcal C|
 \le N-\rho L+|V\setminus Q|+o_D(L),
\]
along with the trivial bound $\log_a|\mathcal C|\le N$.
No price bound on the genie, linearity assumption on the code, or degree
bound on the relation is required.
\end{theorem}

\begin{informalproof}
Apply \cref{thm:rw-deficit} to the cover from \cref{lem:cd-row-words},
with $r=D+2$, $\ell=\ell_{\max}$, and $z=z_D(d)$.
The cover has $K$ independent of depth, and $z_D(d)=o_D(L)$.
Its finite bound is exactly the displayed $L_D(d)$, and its asymptotic
coefficient is $\rho=1/(8K)$. The dimension threshold, unlike $\rho$,
may depend on $D$.
\end{informalproof}

\begin{remark}[What the gap establishes]\label{rem:cd-scope}
The same native trace simultaneously has the robust connected demand pairs
of \cref{thm:cd-robust-demands} and the all-code deficit of
\cref{thm:cd-uniform-gap}.  Its independent rectangle sacrifices a
fixed positive fraction of coordinates.  It forbids internal recovery
edges, and does not claim to separate almost all paired ports into
different components or to eliminate longer paired paths.

The row-word proof depends only on the full head positions, raw-read
positions, and append times, not on the semantic meaning of the shear's
updates.  Thus its relation-level code bound also applies if one changes
the bounded-alphabet local rules while keeping exactly that trace.  A
cheap-genie fibre count would exclude a transposer on this specified
trace: choose price less than $\rho L$, fix the resulting depth, apply
the full-code bound to each genie fibre, and sum the fibre cardinalities.
We do not present that consequence as a new algorithm exclusion.  The
first paper's pass theorem~\cite[Theorem P]{BrendeTapeCoding} already excludes this fixed trace after
absorbing its sublinear nonbank workspace into finite control at each
dimension.  The new content here is the uniform-over-genies, every-induced-
code statement about the actual cone relation.

The fixed number of phases, full rewinds, monotone buffer-length sequences,
small controller span, and common scan-event chart are essential verified
properties of this example.  No normalization of an arbitrary native
schedule to these properties has been proved.  In particular the theorem
does not supply a deficit for all cheap constant-depth native cones, an
adaptive common relation, or the general transposition lower bound.
\end{remark}

\section{What a transposition construction would need}\label{sec:construction}

The native examples distinguish inexpensive execution from inexpensive transposition.
Expansion of the incidence graph and connected demand pairs can occur in linear time,
but neither feature alone defeats a storage-code bound. Conversely, a construction
cannot be excluded by a complex-linear theorem unless its actual genie fibres lie in
the theorem's linear class with the stated coefficient and block-width bounds.

A candidate oblivious linear-time transposer must therefore fit both sides of the
interface: its trace must evade the available cheap depth-reduction theorems, or its
full finite-alphabet recovery relation must support enough codewords to absorb the
cost of every applicable genie. The connected-demand example shows why a positive
coding deficit uniform in fixed depth is useful for this comparison. It does not
reduce an arbitrary candidate to the row interfaces analyzed here.

For an adaptive construction, fixing one trace is insufficient: any comparison must
also account for the information needed to specify that trace and its recovery sets.
The following section states these remaining obligations precisely.

\section{The remaining coding problem}\label{sec:open}

The structural theorem suggests a finite-alphabet question that mentions neither a transition function nor a choice of local encoding gates.

\begin{question}[All-code paired-distance bound]\label{prob:open-paired}
Fix a degree bound~$\Delta$. Let $N_j$ be recovery relations on finite coordinate sets $V_j$ with $|V_j|\to\infty$. Suppose their incidence graphs have maximum degree at most~$\Delta$ and satisfy the paired-distance condition. Must every sequence of nonempty finite-alphabet storage codes $\mathcal C_j\subseteq\mathcal A_j^{V_j}$ satisfy
\[
 \limsup_{j\to\infty}\frac{\log_{|\mathcal A_j|}|\mathcal C_j|}{|V_j|}\le\frac12?
\]
Here $|\mathcal A_j|\ge2$; the alphabets may vary. The fixed binary alphabet is already an important open case.
\end{question}

The directed analytic theorem answers the linear, bounded-complex-coefficient version at fixed vector width. The hyperfinite theorem answers the all-code question under an additional separator hypothesis. The native examples show that this separator hypothesis is not a consequence of a linear-time origin. Their separate coding bounds show, at the same time, that its failure is not evidence for a high-rate code.

\begin{proposition}[What an affirmative answer would supply]\label{prop:open-transfer}
\notformalized
If \cref{prob:open-paired} has an affirmative answer for every fixed~$\Delta$, then for every fixed tape count, linear-time constant, and cone depth, all finite-alphabet storage codes on the corresponding actual transposition-labelled cone relations have rate at most $1/2+o(1)$, uniformly over those schedules and genies.
\end{proposition}
\begin{informalproof}
Let $B$ be the fixed upper bound on recovery-set sizes. For any fixed $0<\delta<1$, condition on the coordinates whose left incidence degree exceeds an integer $M\ge B/\delta$. There are at most $\delta L$ such coordinates. The retained relation has maximum degree at most $\max(B,M)$, and the native structural theorem gives $b_E(r)=O_r(L/\log d)=o(L)$ for every fixed~$r$. Its number of coordinates is at least $(1-\delta)L$, so the same condition holds relative to its size.

An affirmative answer gives half-rate plus $o(1)$ on the retained relation. By \cref{lem:pre-conditioning}, a full code has logarithmic size at most
\[
 |H|+\tfrac12(L-|H|)+o(L)
 \le\tfrac12L+\tfrac12\delta L+o(L).
\]
First take $d\to\infty$ with $\delta$ fixed, then let $\delta\downarrow0$. Uniformity follows by contradiction: a violating sequence of schedules, genies, and codes would itself be a sequence to which \cref{prob:open-paired} applies after the same fixed-threshold trimming.
\end{informalproof}

This implication requires no assertion that the input and output ports are disconnected in the finite graphs. It also does not require that the recovery relation be symmetric. Conversely, a counterexample to \cref{prob:open-paired} would defeat this proposed abstract coding theorem, not construct a transposer: the counterexample might have no native realization.

\subsection{Depth, rate, and the order of choices}
The transposition application still needs a cheap genie. The fourth paper supplies it for home-based head schedules and certain other restricted motion classes, but the general depth-reduction problem remains open~\cite{BrendeCurved}. If the all-code half bound held, any such class admitting genies of an arbitrarily small fixed positive cost would be excluded from oblivious linear-time transposition.

The exact value $1/2$ is not essential to that comparison. A deficit $\rho>0$ works if one can choose a genie of cost $\eps L$ with $\eps<\rho$ and then apply the coding bound at its resulting fixed depth. A bound $\rho(D)>0$ for each~$D$ is not by itself enough: making the genie cheaper can increase~$D$, and the deficit may then fall below its cost. This is why the depth-independent leading constant in the connected-demand example is useful, even though it is smaller than one half.

For the specific native examples, the conclusion concerns all codes on their full cone relations. Their schedules already lie within classes excluded by the first paper's pass bounds. Extending the relation-level theorem to more general schedules would be new work; the examples are not a reduction of arbitrary machines to a finite set of scans.

\subsection{Finite fields and nonlinear recovery}
The complex theorem controls neither modular rank collapse nor nonlinear recovery. Bounded integer lifts expose modular rank loss as torsion; they do not rule it out. The appendix's matrix certificates give an algebraic way to prove finite-field bounds without controlling coefficients or vector width, but only when the required matrix labels exist. Establishing those labels for arbitrary paired-distance systems, or finding a different all-code inequality, would close a genuine remaining interface.

The linear-algebraic and nonlinear statements should therefore be kept separate. A machine's genie is an arbitrary function of its binary input even when its head schedule is fixed. No replacement of that function by a complex linear map is available from the geometry alone.

\subsection{Adaptive schedules}
The geometric incidence theorem applies to each fixed legal adaptive trace. Its proof transports abstract labels along frozen head walks and does not randomize the inputs that select the trace. To use it for an adaptive machine, however, one must control the family of inputs producing that trace and the information needed to identify the relevant recovery relation. Different inputs can have different cones, output times, and genies. The tracewise theorem supplies none of that description for free.

Thus the paper isolates a specific structural invariant and several coding mechanisms that respect it. A complete lower-bound argument still needs a sufficiently strong finite-alphabet capacity bound, a compatible cheap depth reduction, and, for unrestricted~(T), a paid all-input treatment of adaptivity.

\appendix
\section{Finite matrix certificates and an expanding example}
\label{sec:mc-section}

The component phases in \cref{thm:dc-complex} have a finite algebraic
counterpart.  Its witnesses may have growing dimension while the code
being bounded retains any chosen block width.  This distinction permits
a finite-field result without a spectral limit of recovery matrices.
Central layers then transfer its bound to a specified class of nonlinear
finite-group recovery rules.

\begin{lemma}[Rank identity]
\formal{RankCertificate.lean}{rank_identity}
\label{lem:mc-identity}
For endomorphisms $A,L,R$ of any finite-dimensional vector space over any
field $F$,
\begin{equation}
 2\operatorname{rank}(I-A)
 \ge\operatorname{rank}(L-R)-\operatorname{rank}(AL-RA).
 \label{eq:mc-rank-identity}
\end{equation}
No invertibility or commutativity is required.
\end{lemma}
\begin{outline}
For $M=I-A$ the exact identity is
$L-R=(AL-RA)+ML-RM$.  Apply rank subadditivity and
$\operatorname{rank}(ML),\operatorname{rank}(RM)\le\operatorname{rank}M$.
\end{outline}

\begin{proposition}[Matrix-label certificate]
\formal{MatrixCertificate.lean}{matrix_label_certificate,matrix_label_certificate_extension}
\label{prop:mc-certificate}
Let $A$ have $t\times t$ blocks indexed by $N$ coordinates, supported on
a paired incidence relation $E$.  Give its ports arbitrary labels
$X_p,Y_q\in\operatorname{Mat}_m(F)$, where $m\ge1$.  Then
\begin{equation}
 \operatorname{rank}_F(I-A)\ge
 \frac{t\sum_{p}\operatorname{rank}_F(X_p-Y_p)
 -\sum_{L_pR_q\in E}\operatorname{rank}_F(A_{qp})
                         \operatorname{rank}_F(X_p-Y_q)}{2m}.
 \label{eq:mc-certificate-equation}
\end{equation}
Labels over a field extension of $F$ are also valid witnesses, with
the ranks of the original matrices unchanged.
\end{proposition}
\begin{outline}
Use $\widehat A=A\otimes I_m$ and block-diagonal operators
$L_{pp}=I_t\otimes X_p$, $R_{qq}=I_t\otimes Y_q$.
Then
\[
 \operatorname{rank}(I-\widehat A)=m\operatorname{rank}(I-A),\qquad
 \operatorname{rank}(L-R)=t\sum_p\operatorname{rank}(X_p-Y_p),
\]
and the $(q,p)$ block of $\widehat A L-R\widehat A$ is
$A_{qp}\otimes(X_p-Y_q)$.  Tensor-product ranks multiply over any field.
Writing a block matrix as the sum of its single-block inclusions bounds
its rank by the sum of their ranks.  \cref{lem:mc-identity} gives the
displayed inequality. Scalar extension preserves matrix rank, so the same
argument applies to labels over any extension field.
\end{outline}

In particular, full-rank differences on matched ports and low-rank
differences on incidence edges give a useful certificate.  The following
construction shows that this situation is compatible with uniform
expansion and large matched distances.

\begin{theorem}[An existential characteristic-two family]
\notformalized
\label{thm:mc-expanding-family}
There are constants $\Delta\ge1$, $h>0$, and $m_0$ such that for every
integer $m\ge m_0$ there is a simple paired incidence graph $E_m$ on
$2N_m$ ports, where
\[
 N_m=2^{m^2}(2^m-1)^2,
\]
with the following properties:
\begin{enumerate}
\item Its maximum degree is at most $\Delta$, and every port set $S$ with
      $|S|\le N_m$ has at least $h|S|$ edges leaving it.
\item Every matched pair has distance at least $m$.
\item For every field $F$ of characteristic two, every integer $t\ge1$,
      and every recovery matrix $A$ of block width $t$ supported on $E_m$,
\begin{equation}
 \operatorname{rank}_F(I-A)
 \ge\frac{tN_m}{2}
       -\frac1{2m}\sum_{L_pR_q\in E_m}\operatorname{rank}_F(A_{qp})
 \ge\left(\frac12-\frac{\Delta}{2m}\right)tN_m.
 \label{eq:mc-expander-rank}
\end{equation}
\end{enumerate}
The graph is chosen independently of $F,t,A$.  In particular the bound
is uniform in $t$, even if $t$ grows with $m$.
\end{theorem}

\begin{informalproof}
\emph{An outer graph with rank-one steps.}
Put $s=2^m$ and $D_m=(s-1)^2$.  On the additive group
$\operatorname{Mat}_m(\mathbb F_2)$ join $B$ to $B+uv^{\mathsf T}$
for each ordered pair of nonzero vectors $u,v\in\mathbb F_2^m$.
Each nonzero rank-one matrix has exactly one such representation over
$\mathbb F_2$, so this is a simple $D_m$-regular graph $G_m$.
The edge label $(u,v)$ is unchanged by reversal.

For a matrix $C$ of rank $r$, the additive character
$B\mapsto(-1)^{\sum_{i,j}C_{ij}B_{ij}}$ has normalized eigenvalue
\begin{equation}
 \lambda_r=
 \frac{2^{2m-r}-2^{m+1}+1}{(2^m-1)^2}.
 \label{eq:mc-outer-spectrum}
\end{equation}
Indeed, among the $s-1$ nonzero vectors $u$, exactly $2^{m-r}-1$
satisfy $u^{\mathsf T}C=0$.  Summing the character over nonzero $v$
gives $s-1$ for those $u$ and $-1$ for the others.  The resulting sum
is $2^{2m-r}-2s+1$.  These characters form an orthogonal eigenbasis.
For $m\ge2$ and $r\ge1$, the eigenvalues decrease from at most $1/2$
to $-1/(s-1)\ge-1/3$.  Thus all nonconstant eigenvalues have absolute
value at most $1/2$.

\emph{Reducing the degree.}
Fix $k=1024$.  Theorem~1 of Bordenave~\cite{Bordenave} implies that,
for every sufficiently large $D$, there is a simple $k$-regular graph
$H_D$ on $D$ vertices whose normalized nonconstant eigenvalues have
absolute value at most $1/10$.  In fact
$(2\sqrt{1023}+1)/1024<1/10$, and $kD$ is always even.
Choose one such graph on the label set of $G_m$.

Use the zig-zag construction of Reingold, Vadhan, and
Wigderson~\cite[Section~3]{RVW}.  Its vertex set consists of pairs $(B,a)$,
where $a$ is an outer edge label.  A step first changes $a$ along one
edge of $H_{D_m}$, then takes that labelled edge of $G_m$, and finally
changes the new label along one edge of $H_{D_m}$.
The resulting undirected multigraph $Z_m$ has $N_m$ vertices and degree
$\Delta=k^2$.  Every edge changes $B$ by a matrix of rank exactly one.
There are no loops; we permit parallel edges in the spectral accounting.

For completeness, the following loose spectral estimate suffices.
Let $P$ average within each label cloud, let $\mathcal H$ apply the
normalized adjacency of $H_{D_m}$ within clouds, and let $\mathcal R$
be the outer rotation, an orthogonal involution.
The normalized adjacency of $Z_m$ is
$\mathcal H\mathcal R\mathcal H$.
For a vector $f$ orthogonal to the constants, write
$f=f_\parallel+f_\perp$, where $f_\parallel=Pf$, and put
$e=\mathcal Hf_\perp$.  Then $\|e\|\le\|f_\perp\|/10$.
On cloud-constant vectors $P\mathcal RP$ is the outer normalized
adjacency.  Hence
\[
 \bigl|\langle f,\mathcal H\mathcal R\mathcal Hf\rangle\bigr|
 \le\tfrac12\|f_\parallel\|^2
      +\tfrac2{10}\|f_\parallel\|\|f_\perp\|
      +\tfrac1{100}\|f_\perp\|^2
 \le0.71\|f\|^2.
\]
Since the adjacency is self-adjoint, all its nonconstant eigenvalues
have absolute value at most $\lambda=0.71<1$.

\emph{Ports and their prescribed pairing.}
For a vertex $x=(B,a)$ put $\theta(x)=(B+I_m,a)$.
Take the bipartite double cover of $Z_m$, keep the physical left port
at $x$ named $L_x$, and name the physical right port at $\theta(y)$ by
$R_y$.  Finally remove parallel copies of edges to obtain the simple
graph $E_m$.  Thus
\begin{equation}
 L_xR_y\in E_m\quad\Longleftrightarrow\quad
 x\text{ and }\theta(y)\text{ are adjacent in }Z_m.
 \label{eq:mc-shifted-pairing}
\end{equation}
Give these ports the matrix labels $X_x=B_x$ and $Y_y=B_y+I_m$.
Each edge difference has rank one.  Each matched difference is $I_m$
and has rank $m$.  Along any port path the matrix differences telescope;
rank subadditivity therefore forces every matched path to have at least
$m$ edges.  In particular there are no self-demand edges when $m>1$.

The double cover has eigenvalues equal to the positive and negative
eigenvalues of $Z_m$.  Its eigenvalue $1$ is simple, and every other
eigenvalue is at most $\lambda$; the additional eigenvalue $-1$ is harmless
for this upper bound.  For a port set $S$ of size $q\le N_m$, apply
the spectral gap to $1_S-q/(2N_m)$.  It gives at least
$\Delta(1-\lambda)q/2$ crossing edges counted with multiplicity.
Every simple edge has multiplicity at most $\Delta$.  Deduplication
therefore leaves at least $hq$ crossing edges for
$h=(1-\lambda)/2>0$, and degree at most $\Delta$.

\emph{The rank bound.}
Interpret the binary labels $X_x,Y_y$ in any field of characteristic two.
Their ranks are preserved by extension of the prime field.
\cref{prop:mc-certificate} gives the first inequality in
\eqref{eq:mc-expander-rank}.  The second follows from
$\operatorname{rank}A_{yx}\le t$ and $|E_m|\le\Delta N_m$.
\end{informalproof}

\begin{corollary}[Linear codes on the expanding family]
\partformal{LinearRecovery.lean}{linear_recovery_matrix_exists,linear_code_rate_of_matrix_labels}{Construction of the expanding family and its full-rank matched and rank-one edge labels.}
\label{cor:mc-linear-codes}
Over any characteristic-two field, every linear storage space on $E_m$
with $t$ symbols per coordinate has normalized dimension at most
$1/2+\Delta/(2m)$.  This includes linear spaces initially supplied with
nonlinear recovery functions.
\end{corollary}
\begin{informalproof}
On a linear space, determination of a coordinate by its neighbours means
that the kernel of the neighbour projection is contained in the kernel
of the coordinate projection.  The latter therefore factors linearly
through the image of the former, and this linear map extends to the
ambient neighbour space.  Choose these linear recovery maps for all
coordinates to obtain $A$ with $\mathcal C\subseteq\ker(I-A)$, and use
\cref{thm:mc-expanding-family}.
\end{informalproof}

These graphs are not uniformly hyperfinite.  If deletion of $z$ ports
leaves components of size at most a fixed $K<N_m$, expansion summed
over the components gives
$h(2N_m-z)\le\Delta z$.  Thus every such separator has a fixed positive
linear size.  The preceding finite-field result consequently goes beyond
the hyperfinite hypothesis of \cref{cor:dc-hyperfinite}. The next theorem
also transfers it to nonlinear codes with group-expression recovery.

Two scope points are useful.  First, the binary translation $\theta$ is
an involutive automorphism of $Z_m$, so the demand relation in this
particular construction is symmetric.  Its contribution is a
characteristic-two bound uniform in block width, not an example requiring
asymmetric recovery.  Second, the inner graphs were chosen by an existence
theorem.  No efficient uniform generator or native tape realization is
asserted.  These are abstract relations; their degree, expansion, and
paired-port separation do not identify them with all native cones.

\subsection{From vector-linear bounds to group-expression recovery}

Fix a finite loopless recovery relation $E$ on $N>0$ coordinates and a
prime $p$. For a nontrivial finite $p$-group $G$, an \emph{endomorphic
recovery expression} uses permitted parent variables, constants in $G$,
multiplication, inversion, and fixed group endomorphisms $G\to G$.
Expressions are finite, with arbitrary nesting and repeated variables.
Each row has one fixed expression on the whole code. The code itself
can be any nonempty subset of $G^V$ satisfying these equations; it need
not be a subgroup. Arbitrary symbol maps inside an expression are not
included in this language.

\begin{theorem}[Central-layer transfer and the exact benchmark]
\notformalized
\label{thm:mc-group-transfer}
Suppose that for every integer $t\ge1$ and every $t\times t$ block
matrix $A$ over $\F_p$ supported on $E$,
$\dim\ker(I-A)\le\rho t$. Then every endomorphic-expression code
over every finite $p$-group satisfies $|\mathcal C|\le|G|^\rho$.
In particular, define
\[
 \gamma_p(E)=\sup_{t\ge1}\ max_{A\text{ supported on }E}
                    \frac{\dim_{\F_p}\ker(I-A)}{tN}.
\]
The supremum of $N^{-1}\log_{|G|}|\mathcal C|$ over all nontrivial
finite $p$-groups and their nonempty endomorphic-expression codes is
exactly $\gamma_p(E)$, the all-width vector-linear benchmark on $E$.
\end{theorem}
\begin{informalproof}
Use the lower exponent-$p$ central series
$P_1=G$, $P_{j+1}=P_j^p[P_j,G]$, where products denote generated
subgroups. Its terms are fully invariant. Each
$V_j=P_j/P_{j+1}$ is central in $G/P_{j+1}$ and is a vector space
over $\F_p$, of dimension $t_j$. Every endomorphism of $G$ therefore
induces a linear map on $V_j$.

Here is an elementary termination argument. Every nonzero finite
$\F_p$-representation of $G$ has a nonzero fixed vector: its additive
set has size divisible by $p$, every nontrivial orbit has size divisible
by $p$, and zero is fixed, so there are at least $p$ fixed vectors.
Iterate this observation on quotient representations to obtain a full
flag with trivial one-dimensional factors. On the regular module
$\F_p[G]$ of dimension $D=|G|$, every $g-1$ lowers this flag.
Thus its augmentation ideal $I$ satisfies $I^D=0$.
Inductively, $g\in P_j$ implies $g-1\in I^j$: the assertion is
preserved under products and inverses, powers satisfy
$g^p-1=(g-1)^p\in I^{j+1}$, and
\[
 [g,h]-1=g^{-1}h^{-1}
       \bigl((g-1)(h-1)-(h-1)(g-1)\bigr)\in I^{j+1}.
\]
Consequently $P_D=1$ and $\sum_jt_j=\log_p|G|$.

Project the actual code to $(G/P_j)^V$. Fix one of its attained
quotient assignments and choose lifts $r_v\in G/P_{j+1}$.
Every further lift is uniquely $r_va_v$, with $a_v\in V_j$.
Structural induction on an allowed expression $w$ gives
\[
 w((r_va_v)_v)=w((r_v)_v)L_w(a),
\]
where $L_w$ is linear over $\F_p$ and uses only the variables of $w$.
Variables give projections and constants give zero. Products add the
linear maps, inversion negates them, and an endomorphism composes with
its induced map on $V_j$. Centrality permits these operations even
when the base values do not commute. In this language $L_w$ is
independent of the chosen base assignment.

The recovery equations in the fibre therefore have the form
$(I-A_j)a=b$, with $t_j\times t_j$ blocks supported on $E$.
Their solutions form an affine translate of $\ker(I-A_j)$ when
nonempty. The actual code fibre can be a proper subset of this set;
no assertion that all quotient solutions lift is needed.
Multiply the uniform bounds on these successive fibres to get
\[
 \log_p|\mathcal C|\le\sum_j\dim\ker(I-A_j)
                      \le\rho\sum_jt_j=\rho\log_p|G|.
\]
Zero-width layers are omitted. Trivial groups and empty codes satisfy
the cardinality assertion separately.

Taking $\rho=N\gamma_p(E)$ proves the benchmark upper bound without
assuming its supremum is attained. Conversely take the additive group
$G=(\F_p)^t$. Every block linear map is an allowed endomorphism,
and sums of them realize each kernel in the definition of $\gamma_p$.
Thus every rate in that supremum is realized by the allowed language.
Finally any vector-linear storage space, even with initially nonlinear
decoders, admits linear decoders by factoring its parent projections,
as in \cref{cor:mc-linear-codes}. This identifies the benchmark with
the usual vector-linear one.
\end{informalproof}

Central lifting for ordinary group equations is classical; for example,
Shmel'kin's nonsingular-system theorem is stated in
\cite[Introduction]{Mikheenko}. Here fixed endomorphisms require
fully invariant vector layers, and the application needs bounds uniform
in their widths. We do not claim a new general theorem on solvability
of group equations or an unrestricted nonlinear-to-linear reduction.

\begin{corollary}[Group codes and exceptional recovery rows]
\notformalized
\label{cor:mc-group-codes}
On the family $E_m$ of \cref{thm:mc-expanding-family}, let $G$ be any
nontrivial finite $2$-group. If all recovery rows have endomorphic
expressions, then
\[
 \log_{|G|}|\mathcal C|\le
 \frac{N_m}{2}+\frac{|E_m|}{2m}
 \le\left(\frac12+\frac{\Delta}{2m}\right)N_m.
\]
More generally, suppose only rows outside a specified $J\subseteq V$
are required to have such expressions. With $S=V\setminus J$ and
$E_m[S]$ the induced incidence relation,
\[
 \log_{|G|}|\mathcal C|\le
 \frac{N_m+|J|}{2}+\frac{|E_m[S]|}{2m}
 \le\frac{N_m+|J|}{2}+\frac{\Delta(N_m-|J|)}{2m}.
\]
These bounds are uniform in group size and expression length.
For $m>\Delta$, a normalized rate at least $1/2+\delta$ requires
\[
 \frac{|J|}{N_m}\ge
       \frac{2\delta-\Delta/m}{1-\Delta/m}.
\]
In particular a sustained excess $\delta>0$ forces
$\liminf_{m\to\infty}|J|/N_m\ge2\delta$.
\end{corollary}
\begin{informalproof}
Restrict the same binary port labels to both ports of coordinates in
$S$. Their matched differences still have rank $m$ and their edge
differences rank one. \Cref{prop:mc-certificate} gives, over $\F_2$
and at every width $t$,
\[
 \dim\ker(I-A)\le t\left(\frac{|S|}{2}
                                      +\frac{|E_m[S]|}{2m}\right).
\]
Apply \cref{thm:mc-group-transfer}. For $J\ne\varnothing$, first
fix the actual coordinates on $J$ and substitute their values as
constants. Each remaining fibre has exactly the required expression
form on $S$. Summing over at most $|G|^{|J|}$ fibres pays for these
coordinates once. Use $|E_m[S]|\le\Delta|S|$ and rearrange to obtain
the final estimate. The empty $S$ case is the trivial full-code bound.
\end{informalproof}

Thus this particular expanding family obstructs more than linear code
spaces: finite-group nonlinearity in the stated language cannot improve
its asymptotic half-rate bound. A construction would need a different
relation or a quantitatively substantial failure of that language.
No small exceptional set or such expression form has been supplied for
arbitrary native cones, and group operations have not been assigned
unit cost as tape operations.

\section{Index of Lean declarations}\label{sec:leanindex}

Declarations are relative to \texttt{DirectedRecovery}; their axiom checks are generated together with this index (\cref{sec:verification}).

No new results are formalized in the initial draft.

\section{Proof status and verification}\label{sec:verification}

Every numbered result carries an explicit status line. The finite conditioning, fibre, independent-coordinate, rectangle, port-color,
component, separator, and rank bounds, together with the hyperfinite half-rate
corollary, are Lean-verified. The finite counting part of degree trimming is
formalized; its native premise and asymptotic application remain written proofs.
The remaining new results currently have written proofs. The earlier papers are cited at a pinned revision; their verification
claims apply only to their own stated declarations. The geometric, analytic and
native-example arguments here are written proofs, subject to further scrutiny.
The general all-code paired-distance bound, unrestricted cheap depth reduction,
and an all-input treatment of adaptive schedules remain open.

The Lean library uses Lean and Mathlib v4.34.1 and requires the shared
\texttt{lean-machine} package and the fourth paper's library from the same repository. No new axioms,
\texttt{sorry}, \texttt{admit}, or kernel-bypassing decision procedures are permitted.
The only permitted axioms are \texttt{propext}, \texttt{Classical.choice}, and
\texttt{Quot.sound}. The declaration index and \texttt{PaperAxioms.lean} are generated
from the same inline formal-status annotations by \texttt{tools/lean\_index.py}.

From the paper directory, run \texttt{python3 tools/lean\_index.py} to regenerate
these files. In \texttt{lean/}, run \texttt{lake exe cache get}, \texttt{lake build},
and \texttt{lake env lean PaperAxioms.lean}. The entire repository must be present
with its sibling directories. In \texttt{paper/}, run
\texttt{latexmk -pdf main.tex} to compile the manuscript.

Formal verification, where recorded, establishes the precise Lean statements;
it does not certify external novelty or remove hypotheses from their applications.

\bibliographystyle{plainurl}
\bibliography{references}
\end{document}
